\documentclass[aspectratio=169,8pt]{beamer}
\usetheme{Madrid}
\usepackage[utf8]{inputenc}
\usepackage[T1]{fontenc}
\usepackage{amsmath,amssymb}
\usepackage{booktabs}
\usepackage{enumitem}
\setlist[enumerate]{label=\Alph*.,leftmargin=*,itemsep=1pt,topsep=2pt}
\setlist[itemize]{leftmargin=*,itemsep=1pt,topsep=2pt}
\setbeamerfont{block body}{size=\tiny}
\setbeamerfont{block title}{size=\scriptsize}
\setbeamerfont{frametitle}{size=\footnotesize}
\setbeamerfont{normal text}{size=\tiny}
\setbeamertemplate{navigation symbols}{}
\setbeamertemplate{itemize item}{\tiny$\bullet$}
\setbeamertemplate{itemize subitem}{\tiny$\circ$}

\title[GARP RAI]{GARP Risk \& AI --- Supplementary Memory \& Theory Questions}
\subtitle{Topics Not Covered in the Main 150-Question Deck}
\author{Unofficial}
\date{\today}

\begin{document}
	
	\begin{frame}
		\titlepage
	\end{frame}
	
	% =================================================================
	\section{Classical AI and Deep Learning History}
	% =================================================================
	
	\begin{frame}{S1 --- Who Coined the Term ``Artificial Intelligence''?}
		\begin{block}{Question}
			The term ``artificial intelligence'' was coined in 1956 by:
			\begin{enumerate}
				\item Alan Turing
				\item John McCarthy
				\item Marvin Minsky
				\item Claude Shannon
			\end{enumerate}
		\end{block}
		\begin{exampleblock}{Answer: B}\end{exampleblock}
		\begin{block}{Explanation}
			John McCarthy of Dartmouth College coined the term in 1956, along with Marvin Minsky, Nathan Rochester, and Claude Shannon, at the seminal Dartmouth Summer Research Project.
		\end{block}
	\end{frame}
	
	\begin{frame}{S2 --- Logic Theorist Program}
		\begin{block}{Question}
			The Logic Theorist program, designed to generate proofs in propositional logic, was developed by:
			\begin{enumerate}
				\item Newell, Simon, and Shaw
				\item McCarthy, Minsky, and Shannon
				\item McCulloch and Pitts
				\item Rosenblatt
			\end{enumerate}
		\end{block}
		\begin{exampleblock}{Answer: A}\end{exampleblock}
		\begin{block}{Explanation}
			Allen Newell, Herbert Simon, and Cliff Shaw developed Logic Theorist at the Rand Corporation in 1955. It confirmed most theorems in Chapter 2 of Principia Mathematica.
		\end{block}
	\end{frame}
	
	\begin{frame}{S3 --- McCulloch and Pitts (1943)}
		\begin{block}{Question}
			The 1943 paper ``A logical calculus of the ideas immanent in nervous activity'' was authored by:
			\begin{enumerate}
				\item McCulloch and Pitts
				\item Hebb
				\item Rosenblatt
				\item Minsky and Papert
			\end{enumerate}
		\end{block}
		\begin{exampleblock}{Answer: A}\end{exampleblock}
		\begin{block}{Explanation}
			Warren McCulloch and Walter Pitts published this foundational paper in the Bulletin of Mathematical Biology, describing neurons and neural networks in terms of propositional logic.
		\end{block}
	\end{frame}
	
	\begin{frame}{S4 --- Hebb's Neurophysiological Postulate}
		\begin{block}{Question}
			Donald Hebb's 1949 book that proposed associative learning in neural networks is titled:
			\begin{enumerate}
				\item The Organisation of Behaviour
				\item Perceptrons
				\item Parallel Distributed Processing
				\item Computing Machinery and Intelligence
			\end{enumerate}
		\end{block}
		\begin{exampleblock}{Answer: A}\end{exampleblock}
		\begin{block}{Explanation}
			Hebb's ``The Organisation of Behaviour'' (1949) proposed the neurophysiological postulate that when one neuron repeatedly fires another, the connection strengthens.
		\end{block}
	\end{frame}
	
	\begin{frame}{S5 --- Rosenblatt's Perceptron}
		\begin{block}{Question}
			Frank Rosenblatt's Perceptron was formulated in terms of:
			\begin{enumerate}
				\item Symbolic logic
				\item Probability theory
				\item Propositional calculus
				\item Predicate logic
			\end{enumerate}
		\end{block}
		\begin{exampleblock}{Answer: B}\end{exampleblock}
		\begin{block}{Explanation}
			Rosenblatt (1958) formulated the Perceptron ``in terms of probability theory rather than symbolic logic,'' rejecting symbolic or algorithmic approaches.
		\end{block}
	\end{frame}
	
	\begin{frame}{S6 --- Minsky and Papert's Critique}
		\begin{block}{Question}
			Minsky and Papert's 1969 book that severely undermined enthusiasm for neural networks was titled:
			\begin{enumerate}
				\item Perceptrons
				\item The Organisation of Behaviour
				\item Parallel Distributed Processing
				\item Superintelligence
			\end{enumerate}
		\end{block}
		\begin{exampleblock}{Answer: A}\end{exampleblock}
		\begin{block}{Explanation}
			Minsky and Papert's ``Perceptrons'' (1969) showed that single-layer Perceptrons could not learn simple logical functions like XOR.
		\end{block}
	\end{frame}
	
	\begin{frame}{S7 --- Backpropagation Resurgence}
		\begin{block}{Question}
			The 1986 two-volume work that displaced pessimism about neural networks and encouraged connectionism was:
			\begin{enumerate}
				\item Parallel Distributed Processing
				\item Perceptrons
				\item The Organisation of Behaviour
				\item Computing Machinery and Intelligence
			\end{enumerate}
		\end{block}
		\begin{exampleblock}{Answer: A}\end{exampleblock}
		\begin{block}{Explanation}
			Rumelhart, McClelland, and the PDP Research Group published ``Parallel Distributed Processing'' in 1986, demonstrating how layers of interacting neurons could learn complex functions.
		\end{block}
	\end{frame}
	
	\begin{frame}{S8 --- LeNet-5 (1998)}
		\begin{block}{Question}
			LeNet-5, a convolutional neural network designed for recognition of digits, was developed by:
			\begin{enumerate}
				\item Yann LeCun and colleagues
				\item Krizhevsky, Sutskever, and Hinton
				\item DeepMind
				\item OpenAI
			\end{enumerate}
		\end{block}
		\begin{exampleblock}{Answer: A}\end{exampleblock}
		\begin{block}{Explanation}
			Yann LeCun and colleagues published LeNet-5 in 1998, consisting of 8 layers with convolutional and pooling layers.
		\end{block}
	\end{frame}
	
	\begin{frame}{S9 --- AlexNet (2012)}
		\begin{block}{Question}
			AlexNet won the 2012 ImageNet Large Scale Visual Recognition Challenge with an error rate of:
			\begin{enumerate}
				\item 25.8\%
				\item 15.3\%
				\item 10.0\%
				\item 5.0\%
			\end{enumerate}
		\end{block}
		\begin{exampleblock}{Answer: B}\end{exampleblock}
		\begin{block}{Explanation}
			AlexNet (Krizhevsky, Sutskever, Hinton) achieved a 15.3\% error rate in 2012, compared to the previous year's winning 25.8\%.
		\end{block}
	\end{frame}
	
	\begin{frame}{S10 --- AlphaGo vs Lee Sedol}
		\begin{block}{Question}
			In 2016, AlphaGo defeated World Go Champion Lee Sedol with a score of:
			\begin{enumerate}
				\item 5-0
				\item 4-1
				\item 3-2
				\item 2-3
			\end{enumerate}
		\end{block}
		\begin{exampleblock}{Answer: B}\end{exampleblock}
		\begin{block}{Explanation}
			AlphaGo defeated Lee Sedol 4-1 in 2016. It had previously defeated European champion Fan Hui 5-0 in 2015.
		\end{block}
	\end{frame}
	
	\begin{frame}{S11 --- AlphaZero (2017)}
		\begin{block}{Question}
			AlphaZero learned to play chess, Go, and Shogi by:
			\begin{enumerate}
				\item Learning from human game databases
				\item Competing against itself
				\item Following expert-designed heuristics
				\item Using exhaustive search
			\end{enumerate}
		\end{block}
		\begin{exampleblock}{Answer: B}\end{exampleblock}
		\begin{block}{Explanation}
			AlphaZero learned by self-play (competing against itself) without any input from human experts or game databases, defeating Stockfish after only a few hours.
		\end{block}
	\end{frame}
	
	\begin{frame}{S12 --- Stochastic Parrot Term}
		\begin{block}{Question}
			The term ``stochastic parrot'' was coined by:
			\begin{enumerate}
				\item Emily M. Bender
				\item Geoffrey Hinton
				\item Yann LeCun
				\item Fei Fei Li
			\end{enumerate}
		\end{block}
		\begin{exampleblock}{Answer: A}\end{exampleblock}
		\begin{block}{Explanation}
			Emily M. Bender coined the term ``stochastic parrot'' in 2021 to describe LLMs that imitate language without understanding meaning.
		\end{block}
	\end{frame}
	
	\begin{frame}{S13 --- Turing Test Original Name}
		\begin{block}{Question}
			Alan Turing's 1950 paper ``Computing Machinery and Intelligence'' introduced the Turing Test under what original name?
			\begin{enumerate}
				\item The Imitation Game
				\item The Reasoning Test
				\item The Intelligence Game
				\item The Machine Test
			\end{enumerate}
		\end{block}
		\begin{exampleblock}{Answer: A}\end{exampleblock}
		\begin{block}{Explanation}
			Turing called it the ``imitation game,'' where a computer proves intelligence if its textual conversation is indistinguishable from a human's.
		\end{block}
	\end{frame}
	
	\begin{frame}{S14 --- Turing Machine (1936)}
		\begin{block}{Question}
			Alan Turing's 1936 paper introducing the ``Turing machine'' was titled:
			\begin{enumerate}
				\item On Computable Numbers, with an Application to the Entscheidungsproblem
				\item Computing Machinery and Intelligence
				\item The Organisation of Behaviour
				\item A Logical Calculus of the Ideas Immanent in Nervous Activity
			\end{enumerate}
		\end{block}
		\begin{exampleblock}{Answer: A}\end{exampleblock}
		\begin{block}{Explanation}
			Turing's seminal 1936 paper appeared in the Proceedings of the London Mathematical Society and introduced the theoretical computing machine.
		\end{block}
	\end{frame}
	
	% =================================================================
	\section{Unsupervised Learning (Additional Topics)}
	% =================================================================
	
	\begin{frame}{S15 --- K-means as Lloyd's Algorithm}
		\begin{block}{Question}
			The K-means algorithm is also known as:
			\begin{enumerate}
				\item Lloyd's algorithm
				\item Newton's algorithm
				\item Bayes' algorithm
				\item Fisher's algorithm
			\end{enumerate}
		\end{block}
		\begin{exampleblock}{Answer: A}\end{exampleblock}
		\begin{block}{Explanation}
			K-means is also known as Lloyd's algorithm, after Stuart Lloyd who developed it at Bell Labs in 1957.
		\end{block}
	\end{frame}
	
	\begin{frame}{S16 --- Rule of Thumb for K}
		\begin{block}{Question}
			A rough guideline for initially estimating K in K-means is:
			\begin{enumerate}
				\item K = N/2
				\item K = √(N/2)
				\item K = log(N)
				\item K = 2N
			\end{enumerate}
		\end{block}
		\begin{exampleblock}{Answer: B}\end{exampleblock}
		\begin{block}{Explanation}
			A rough guideline is K = √(N/2), where N is the number of data points. This is arbitrary and more sophisticated methods are preferred.
		\end{block}
	\end{frame}
	
	\begin{frame}{S17 --- K-means++ Purpose}
		\begin{block}{Question}
			K-means++ is an improved initialization method whose purpose is to:
			\begin{enumerate}
				\item Reduce K
				\item Spread initial centroids far apart to reduce poor local optima
				\item Increase WCSS
				\item Eliminate outliers
			\end{enumerate}
		\end{block}
		\begin{exampleblock}{Answer: B}\end{exampleblock}
		\begin{block}{Explanation}
			K-means++ selects initial centroids with probability proportional to squared distance from nearest existing centroid, spreading them out and reducing poor local optima.
		\end{block}
	\end{frame}
	
	\begin{frame}{S18 --- Between-Cluster Sum of Squares (BCSS)}
		\begin{block}{Question}
			BCSS measures:
			\begin{enumerate}
				\item Intra-cluster compactness
				\item Sum of squared distances of each cluster's centroid to the overall centroid, weighted by cluster size
				\item Total variance
				\item Number of clusters
			\end{enumerate}
		\end{block}
		\begin{exampleblock}{Answer: B}\end{exampleblock}
		\begin{block}{Explanation}
			BCSS focuses on inter-cluster dissimilarity (separation). A higher BCSS suggests clusters are well separated.
		\end{block}
	\end{frame}
	
	\begin{frame}{S19 --- Hierarchical Clustering: Single Linkage}
		\begin{block}{Question}
			Single linkage defines distance between two clusters as:
			\begin{enumerate}
				\item The greatest distance between any two points in different clusters
				\item The shortest distance between any two points in different clusters
				\item The average distance between all points
				\item The distance between centroids
			\end{enumerate}
		\end{block}
		\begin{exampleblock}{Answer: B}\end{exampleblock}
		\begin{block}{Explanation}
			Single linkage = minimum distance between clusters. Complete linkage = maximum distance. Average linkage = average distance.
		\end{block}
	\end{frame}
	
	\begin{frame}{S20 --- Hierarchical Clustering: Complete Linkage}
		\begin{block}{Question}
			Complete linkage defines distance between two clusters as:
			\begin{enumerate}
				\item The shortest distance between any two points
				\item The greatest distance between any two points in different clusters
				\item The average distance
				\item The distance between medoids
			\end{enumerate}
		\end{block}
		\begin{exampleblock}{Answer: B}\end{exampleblock}
		\begin{block}{Explanation}
			Complete linkage uses the maximum distance between clusters. It tends to produce compact, spherical clusters.
		\end{block}
	\end{frame}
	
	\begin{frame}{S21 --- Divisive vs Agglomerative Clustering}
		\begin{block}{Question}
			Agglomerative clustering starts with:
			\begin{enumerate}
				\item All data in one cluster and splits
				\item Each data point as its own cluster and merges
				\item Random clusters
				\item K clusters
			\end{enumerate}
		\end{block}
		\begin{exampleblock}{Answer: B}\end{exampleblock}
		\begin{block}{Explanation}
			Agglomerative = bottom-up (each point starts as its own cluster, then merge). Divisive = top-down (start with one cluster, then split).
		\end{block}
	\end{frame}
	
	\begin{frame}{S22 --- DBSCAN Core Points}
		\begin{block}{Question}
			In DBSCAN, a core point is defined as:
			\begin{enumerate}
				\item Any point in the dataset
				\item A point with at least a minimum number of other points within a specified radius
				\item A point on the cluster boundary
				\item An outlier
			\end{enumerate}
		\end{block}
		\begin{exampleblock}{Answer: B}\end{exampleblock}
		\begin{block}{Explanation}
			A core point has at least minPts neighbors within radius e. Border points are within ε of a core point. Noise points are neither.
		\end{block}
	\end{frame}
	
	\begin{frame}{S23 --- DBSCAN Noise Points}
		\begin{block}{Question}
			In DBSCAN, a noise point is:
			\begin{enumerate}
				\item A core point
				\item A border point
				\item A point that is neither a core point nor a border point
				\item A centroid
			\end{enumerate}
		\end{block}
		\begin{exampleblock}{Answer: C}\end{exampleblock}
		\begin{block}{Explanation}
			Noise points (outliers) are in low-density regions and are not assigned to any cluster in DBSCAN.
		\end{block}
	\end{frame}
	
	% =================================================================
	\section{Supervised Learning (Additional Topics)}
	% =================================================================
	
	\begin{frame}{S24 --- Gini Coefficient Formula}
		\begin{block}{Question}
			The Gini coefficient for a node with J classes is:
			\begin{enumerate}
				\item $1 - \sum_{j=1}^{J} p_j^2$
				\item $-\sum_{j=1}^{J} p_j \log_2(p_j)$
				\item $\sum_{j=1}^{J} p_j$
				\item $\max_j p_j$
			\end{enumerate}
		\end{block}
		\begin{exampleblock}{Answer: A}\end{exampleblock}
		\begin{block}{Explanation}
			Gini = 1 - summaton pj². Entropy =sumamton pjlog2(pj). Both measure node impurity in classification trees.
		\end{block}
	\end{frame}
	
	\begin{frame}{S25 --- Entropy Formula}
		\begin{block}{Question}
			Entropy for a node with J classes is:
			\begin{enumerate}
				\item $1 - \sum p_j^2$
				\item $-\sum_{j=1}^{J} p_j \log_2(p_j)$
				\item $\sum p_j$
				\item $\max p_j$
			\end{enumerate}
		\end{block}
		\begin{exampleblock}{Answer: B}\end{exampleblock}
		\begin{block}{Explanation}
			Entropy   It measures disorder. Base-2 log is used; changing base scales all values by a constant.
		\end{block}
	\end{frame}
	
	\begin{frame}{S26 --- Pruning: Pre-pruning vs Post-pruning}
		\begin{block}{Question}
			Pre-pruning (online pruning) in decision trees:
			\begin{enumerate}
				\item Grows the tree fully then removes branches
				\item Stops tree growth early using stopping criteria
				\item Removes outliers
				\item Increases tree depth
			\end{enumerate}
		\end{block}
		\begin{exampleblock}{Answer: B}\end{exampleblock}
		\begin{block}{Explanation}
			Pre-pruning prevents a tree from growing too much (e.g., minimum observations per node). Post-pruning grows fully then removes weak branches.
		\end{block}
	\end{frame}
	
	\begin{frame}{S27 --- Reduced Error Pruning}
		\begin{block}{Question}
			Reduced error pruning is a:
			\begin{enumerate}
				\item Top-down approach
				\item Bottom-up approach that replaces nodes with most popular class if validation performance doesn't worsen
				\item Type of boosting
				\item Type of bagging
			\end{enumerate}
		\end{block}
		\begin{exampleblock}{Answer: B}\end{exampleblock}
		\begin{block}{Explanation}
			Reduced error pruning is bottom-up: start at leaves, replace nodes with their most popular class if the pruned tree performs no worse on validation data.
		\end{block}
	\end{frame}
	
	\begin{frame}{S28 --- Cost Complexity Pruning}
		\begin{block}{Question}
			Cost complexity pruning adds a penalty term to RSS, with the trade-off controlled by:
			\begin{enumerate}
				\item α (alpha)
				\item λ (lambda)
				\item η (eta)
				\item μ (mu)
			\end{enumerate}
		\end{block}
		\begin{exampleblock}{Answer: A}\end{exampleblock}
		\begin{block}{Explanation}
			Cost complexity pruning uses parameter a, chosen by cross-validation, to trade off accuracy vs number of terminal nodes.
		\end{block}
	\end{frame}
	
	\begin{frame}{S29 --- Random Forest Feature Subsetting}
		\begin{block}{Question}
			In a random forest, the number of features considered at each split is typically:
			\begin{enumerate}
				\item All features
				\item Approximately the square root of the total number of features
				\item One feature
				\item Half the features
			\end{enumerate}
		\end{block}
		\begin{exampleblock}{Answer: B}\end{exampleblock}
		\begin{block}{Explanation}
			Random forest uses a random subset of features at each split, usually = rootof (total features). This reduces correlation across trees.
		\end{block}
	\end{frame}
	
	\begin{frame}{S30 --- Bagging vs Pasting}
		\begin{block}{Question}
			Pasting differs from bagging because it:
			\begin{enumerate}
				\item Samples with replacement
				\item Samples without replacement
				\item Uses all data
				\item Uses no data
			\end{enumerate}
		\end{block}
		\begin{exampleblock}{Answer: B}\end{exampleblock}
		\begin{block}{Explanation}
			Bagging = sampling with replacement. Pasting = sampling without replacement. Both create multiple subsamples for ensemble training.
		\end{block}
	\end{frame}
	
	\begin{frame}{S31 --- KNN as Lazy Learner}
		\begin{block}{Question}
			KNN is called a ``lazy learner'' because:
			\begin{enumerate}
				\item It trains very slowly
				\item It doesn't build an explicit model during training; computation happens at prediction time
				\item It uses few features
				\item It ignores training data
			\end{enumerate}
		\end{block}
		\begin{exampleblock}{Answer: B}\end{exampleblock}
		\begin{block}{Explanation}
			KNN memorizes the training dataset and performs all computation at prediction time when distances are calculated.
		\end{block}
	\end{frame}
	
	\begin{frame}{S32 --- KNN K Selection}
		\begin{block}{Question}
			A common choice for K in KNN is:
			\begin{enumerate}
				\item K = N
				\item K = √N
				\item K = 1
				\item K = N/2
			\end{enumerate}
		\end{block}
		\begin{exampleblock}{Answer: B}\end{exampleblock}
		\begin{block}{Explanation}
			K = √N is a common heuristic. Small K overfits; large K underfits. Cross-validation can tune K.
		\end{block}
	\end{frame}
	
	\begin{frame}{S33 --- SVM Support Vectors}
		\begin{block}{Question}
			Support vectors are:
			\begin{enumerate}
				\item All training points
				\item The training points closest to the decision boundary that define the margin
				\item The centroids of classes
				\item Outliers
			\end{enumerate}
		\end{block}
		\begin{exampleblock}{Answer: B}\end{exampleblock}
		\begin{block}{Explanation}
			Support vectors lie on the margin edges and determine the position and orientation of the decision boundary.
		\end{block}
	\end{frame}
	
	\begin{frame}{S34 --- SVM Maximum Margin Classifier}
		\begin{block}{Question}
			The maximum margin classifier is:
			\begin{enumerate}
				\item Any separating line
				\item The line equidistant from margin constraints that maximizes the margin
				\item The line through the origin
				\item The line minimizing margin
			\end{enumerate}
		\end{block}
		\begin{exampleblock}{Answer: B}\end{exampleblock}
		\begin{block}{Explanation}
			The maximum margin classifier is the optimal decision boundary that maximizes the margin between classes.
		\end{block}
	\end{frame}
	
	\begin{frame}{S35 --- SVM Soft Margin Parameter}
		\begin{block}{Question}
			In SVM soft margin, a large λ (lambda) leads to:
			\begin{enumerate}
				\item Narrower margin, fewer misclassifications
				\item Wider margin, more misclassifications allowed
				\item No margin
				\item Perfect separation
			\end{enumerate}
		\end{block}
		\begin{exampleblock}{Answer: B}\end{exampleblock}
		\begin{block}{Explanation}
			Large lambda allows wider margin with more misclassifications (more tolerant). Small lambda penalizes misclassification heavily (narrower margin).
		\end{block}
	\end{frame}
	
	\begin{frame}{S36 --- SVM Kernel Trick Purpose}
		\begin{block}{Question}
			The kernel trick in SVM allows:
			\begin{enumerate}
				\item Linear classification only
				\item Nonlinear decision boundaries by projecting data to higher-dimensional space without explicit computation
				\item Removing outliers
				\item Reducing features
			\end{enumerate}
		\end{block}
		\begin{exampleblock}{Answer: B}\end{exampleblock}
		\begin{block}{Explanation}
			The kernel trick computes inner products in higher-dimensional space without explicitly transforming data, enabling nonlinear boundaries.
		\end{block}
	\end{frame}
	
	\begin{frame}{S37 --- Universal Approximation Theorem}
		\begin{block}{Question}
			The universal approximation theorem states that:
			\begin{enumerate}
				\item Neural networks can only approximate linear functions
				\item Neural networks with enough hidden neurons can approximate any function with arbitrary precision
				\item Neural networks cannot learn nonlinear functions
				\item Neural networks always overfit
			\end{enumerate}
		\end{block}
		\begin{exampleblock}{Answer: B}\end{exampleblock}
		\begin{block}{Explanation}
			The theorem states that with the right number of hidden neurons, neural networks can approximate any function with arbitrary precision under some assumptions.
		\end{block}
	\end{frame}
	
	\begin{frame}{S38 --- Sigmoid Activation Function}
		\begin{block}{Question}
			The sigmoid (logistic) activation function is:
			\begin{enumerate}
				\item $f(z) = \max(0, z)$
				\item $f(z) = 1/(1 + e^{-z})$
				\item $f(z) = (e^{2z} - 1)/(e^{2z} + 1)$
				\item $f(z) = z$
			\end{enumerate}
		\end{block}
		\begin{exampleblock}{Answer: B}\end{exampleblock}
		\begin{block}{Explanation}
			Sigmoid = 1/ soemthign , outputs between 0 and 1. Used in binary classification output layers.
		\end{block}
	\end{frame}
	
	\begin{frame}{S39 --- ReLU Activation Function}
		\begin{block}{Question}
			The ReLU activation function is:
			\begin{enumerate}
				\item $f(z) = 1/(1 + e^{-z})$
				\item $f(z) = \max(0, z)$
				\item $f(z) = \tanh(z)$
				\item $f(z) = z^2$
			\end{enumerate}
		\end{block}
		\begin{exampleblock}{Answer: B}\end{exampleblock}
		\begin{block}{Explanation}
			ReLU = max(0, z). It is computationally efficient and less prone to vanishing gradient for positive values.
		\end{block}
	\end{frame}
	
	\begin{frame}{S40 --- Leaky ReLU}
		\begin{block}{Question}
			The leaky ReLU activation function is:
			\begin{enumerate}
				\item $f(z) = \max(0, z)$
				\item $f(z) = 0.01z$ if $z < 0$, $z$ if $z \geq 0$
				\item $f(z) = 1/(1+e^{-z})$
				\item $f(z) = \tanh(z)$
			\end{enumerate}
		\end{block}
		\begin{exampleblock}{Answer: B}\end{exampleblock}
		\begin{block}{Explanation}
			Leaky ReLU allows small negative values (0.01z), preventing dead neurons that can occur with standard ReLU.
		\end{block}
	\end{frame}
	
	\begin{frame}{S41 --- Hyperbolic Tangent Activation}
		\begin{block}{Question}
			The hyperbolic tangent (tanh) activation function outputs values between:
			\begin{enumerate}
				\item 0 and 1
				\item -1 and 1
				\item 0 and 
				\item -infinity and infinity
			\end{enumerate}
		\end{block}
		\begin{exampleblock}{Answer: B}\end{exampleblock}
		\begin{block}{Explanation}
			Tanh outputs between -1 and 1, similar in shape to sigmoid but zero-centered. .
		\end{block}
	\end{frame}
	
	\begin{frame}{S42 --- Dropout Regularization}
		\begin{block}{Question}
			Dropout regularization works by:
			\begin{enumerate}
				\item Removing features
				\item Randomly deactivating a fraction of neurons during training
				\item Increasing learning rate
				\item Adding more layers
			\end{enumerate}
		\end{block}
		\begin{exampleblock}{Answer: B}\end{exampleblock}
		\begin{block}{Explanation}
			Dropout randomly drops neurons during training, preventing co-adaptation and reducing overfitting. At inference, all neurons are used.
		\end{block}
	\end{frame}
	
	\begin{frame}{S43 --- Early Stopping}
		\begin{block}{Question}
			Early stopping in neural network training:
			\begin{enumerate}
				\item Stops training when validation error starts to rise
				\item Stops after a fixed number of epochs
				\item Never stops
				\item Stops immediately
			\end{enumerate}
		\end{block}
		\begin{exampleblock}{Answer: A}\end{exampleblock}
		\begin{block}{Explanation}
			Early stopping monitors validation performance and halts training when it starts to deteriorate, preventing overfitting.
		\end{block}
	\end{frame}
	
	\begin{frame}{S44 --- Batch Normalization}
		\begin{block}{Question}
			Batch normalization:
			\begin{enumerate}
				\item Normalizes weights
				\item Adds normalization layers between hidden layers to stabilize training
				\item Removes features
				\item Increases learning rate
			\end{enumerate}
		\end{block}
		\begin{exampleblock}{Answer: B}\end{exampleblock}
		\begin{block}{Explanation}
			Batch normalization normalizes inputs to each layer across the batch, stabilizing training and mitigating internal covariate shift.
		\end{block}
	\end{frame}
	
	\begin{frame}{S45 --- CNN Convolutional Layer}
		\begin{block}{Question}
			A convolutional layer in a CNN:
			\begin{enumerate}
				\item Connects every neuron to every other neuron
				\item Applies a filter/kernel to local regions of the input
				\item Only works with 1D data
				\item Removes features
			\end{enumerate}
		\end{block}
		\begin{exampleblock}{Answer: B}\end{exampleblock}
		\begin{block}{Explanation}
			Convolutional layers use filters that slide across the input, creating feature maps. This reduces parameters and builds translation invariance.
		\end{block}
	\end{frame}
	
	\begin{frame}{S46 --- CNN Pooling Layer}
		\begin{block}{Question}
			A pooling layer in a CNN:
			\begin{enumerate}
				\item Increases dimensions
				\item Reduces dimensions by summarizing regions (e.g., average or max)
				\item Adds parameters
				\item Removes activation functions
			\end{enumerate}
		\end{block}
		\begin{exampleblock}{Answer: B}\end{exampleblock}
		\begin{block}{Explanation}
			Pooling layers reduce spatial dimensions by taking summary statistics (e.g., max, average) over regions, reducing computation.
		\end{block}
	\end{frame}
	
	\begin{frame}{S47 --- Autoencoder Purpose}
		\begin{block}{Question}
			Autoencoders are primarily used for:
			\begin{enumerate}
				\item Classification
				\item Dimensionality reduction and anomaly detection
				\item Regression
				\item Clustering only
			\end{enumerate}
		\end{block}
		\begin{exampleblock}{Answer: B}\end{exampleblock}
		\begin{block}{Explanation}
			Autoencoders learn compressed representations (dimensionality reduction) and can detect anomalies via high reconstruction error.
		\end{block}
	\end{frame}
	
	\begin{frame}{S48 --- Autoencoder vs PCA}
		\begin{block}{Question}
			Autoencoders differ from PCA because they:
			\begin{enumerate}
				\item Only perform linear transformations
				\item Can capture nonlinear relationships using activation functions
				\item Require no training
				\item Always produce orthogonal components
			\end{enumerate}
		\end{block}
		\begin{exampleblock}{Answer: B}\end{exampleblock}
		\begin{block}{Explanation}
			Autoencoders with nonlinear activation functions can capture nonlinear patterns, unlike PCA which is strictly linear.
		\end{block}
	\end{frame}
	
	\begin{frame}{S49 --- Autoencoder Loss Function}
		\begin{block}{Question}
			The autoencoder loss function is:
			\begin{enumerate}
				\item $\sum \sum (x_{ij} - \hat{x}_{ij})^2$
				\item $\sum |x_{ij} - \hat{x}_{ij}|$
				\item $\max(x_{ij} - \hat{x}_{ij})$
				\item $\sum x_{ij}$
			\end{enumerate}
		\end{block}
		\begin{exampleblock}{Answer: A}\end{exampleblock}
		\begin{block}{Explanation}
			Autoencoder loss = reconstruction error . The goal is to minimize this error.
		\end{block}
	\end{frame}
	
	% =================================================================
	\section{Model Estimation and Evaluation (Additional Topics)}
	% =================================================================
	
	\begin{frame}{S50 --- Ordinary Least Squares (OLS)}
		\begin{block}{Question}
			OLS estimates parameters by minimizing:
			\begin{enumerate}
				\item Sum of residuals
				\item Residual Sum of Squares (RSS)
				\item Mean Absolute Error
				\item Maximum likelihood
			\end{enumerate}
		\end{block}
		\begin{exampleblock}{Answer: B}\end{exampleblock}
		\begin{block}{Explanation}
			OLS minimizes RSS  . Squaring prevents positive and negative residuals from canceling.
		\end{block}
	\end{frame}
	
	\begin{frame}{S51 --- OLS Slope Formula}
		\begin{block}{Question}
			The OLS slope estimate is:
			\begin{enumerate}
				\item $\hat{\beta}_1 = \bar{y} - \hat{\beta}_0 \bar{x}$
				\item $\hat{\beta}_1 = \dfrac{\sum y_i x_i - N\bar{y}\bar{x}}{\sum x_i^2 - N\bar{x}^2}$
				\item $\hat{\beta}_1 = \bar{y}/\bar{x}$
				\item $\hat{\beta}_1 = \sum(y_i - \bar{y})(x_i - \bar{x})$
			\end{enumerate}
		\end{block}
		\begin{exampleblock}{Answer: B}\end{exampleblock}
		\begin{block}{Explanation}
			OLS slope = covariance(x,y)/variance(x). Option A is the intercept formula. Option D is missing the denominator.
		\end{block}
	\end{frame}
	
	\begin{frame}{S52 --- OLS Intercept Formula}
		\begin{block}{Question}
			Given $\hat{\beta}_1$, the OLS intercept is:
			\begin{enumerate}
				\item $\hat{\beta}_0 = \bar{y} - \hat{\beta}_1 \bar{x}$
				\item $\hat{\beta}_0 = \bar{y} + \hat{\beta}_1 \bar{x}$
				\item $\hat{\beta}_0 = \bar{y}/\hat{\beta}_1 \bar{x}$
				\item $\hat{\beta}_0 = \hat{\beta}_1 \bar{x}$
			\end{enumerate}
		\end{block}
		\begin{exampleblock}{Answer: A}\end{exampleblock}
		\begin{block}{Explanation}
			The regression line passes through the point of means
		\end{block}
	\end{frame}
	
	\begin{frame}{S53 --- Maximum Likelihood Estimation}
		\begin{block}{Question}
			MLE estimates parameters by:
			\begin{enumerate}
				\item Minimizing RSS
				\item Maximizing the likelihood of observing the data
				\item Minimizing absolute errors
				\item Maximizing R²
			\end{enumerate}
		\end{block}
		\begin{exampleblock}{Answer: B}\end{exampleblock}
		\begin{block}{Explanation}
			MLE finds parameter values that maximize the likelihood function L = P(X|θ). For linear regression with normal errors, MLE = OLS.
		\end{block}
	\end{frame}
	
	\begin{frame}{S54 --- Log-Likelihood Function (Logit)}
		\begin{block}{Question}
			The log-likelihood for the logit model is:
			\begin{enumerate}
				\item $\sum [y_i \ln F(y_i) + (1-y_i)\ln(1-F(y_i))]$
				\item $\sum y_i$
				\item $\sum (y_i - F(y_i))^2$
				\item $\prod F(y_i)$
			\end{enumerate}
		\end{block}
		\begin{exampleblock}{Answer: A}\end{exampleblock}
		\begin{block}{Explanation}
			Log-likelihood sums log-probabilities. Option D is the likelihood before log. Option C is squared-error (not likelihood).
		\end{block}
	\end{frame}
	
	\begin{frame}{S55 --- Gradient Descent Update Rule}
		\begin{block}{Question}
			The gradient descent weight update is:
			\begin{enumerate}
				\item $w_i^{\text{new}} = w_i^{\text{old}} + \eta \frac{\partial L}{\partial w_i}$
				\item $w_i^{\text{new}} = w_i^{\text{old}} - \eta \frac{\partial L}{\partial w_i}$
				\item $w_i^{\text{new}} = w_i^{\text{old}} \cdot \eta$
				\item $w_i^{\text{new}} = \eta$
			\end{enumerate}
		\end{block}
		\begin{exampleblock}{Answer: B}\end{exampleblock}
		\begin{block}{Explanation}
			Gradient descent subtracts the learning-rate-scaled gradient. Option A would be gradient ascent.
		\end{block}
	\end{frame}
	
	\begin{frame}{S56 --- Learning Rate Decay (Exponential)}
		\begin{block}{Question}
			Exponential decay of learning rate is:
			\begin{enumerate}
				\item $\eta_t = \eta_0 + kt$
				\item $\eta_t = \eta_0 e^{-kt}$
				\item $\eta_t = \eta_0 \cdot kt$
				\item $\eta_t = \eta_0/k$
			\end{enumerate}
		\end{block}
		\begin{exampleblock}{Answer: B}\end{exampleblock}
		\begin{block}{Explanation}
			Exponential decay   Inverse decay   
		\end{block}
	\end{frame}
	
	\begin{frame}{S57 --- Learning Rate Decay (Inverse)}
		\begin{block}{Question}
			Inverse decay of learning rate is:
			\begin{enumerate}
				\item $\eta_t = \eta_0 e^{-kt}$
				\item $\eta_t = \eta_0/(1+kt)$
				\item $\eta_t = \eta_0(1+kt)$
				\item $\eta_t = \eta_0 \cdot k$
			\end{enumerate}
		\end{block}
		\begin{exampleblock}{Answer: B}\end{exampleblock}
		\begin{block}{Explanation}
			Inverse decay . Both exponential and inverse decay reduce learning rate over epochs.
		\end{block}
	\end{frame}
	
	\begin{frame}{S58 --- Backpropagation Purpose}
		\begin{block}{Question}
			Backpropagation is used to:
			\begin{enumerate}
				\item Forward propagate inputs
				\item Compute gradients efficiently using chain rule from output to input layers
				\item Cluster data
				\item Reduce features
			\end{enumerate}
		\end{block}
		\begin{exampleblock}{Answer: B}\end{exampleblock}
		\begin{block}{Explanation}
			Backpropagation computes gradients layer-by-layer using the chain rule, starting from the output layer and working backward.
		\end{block}
	\end{frame}
	
	\begin{frame}{S59 --- Vanishing Gradient Problem}
		\begin{block}{Question}
			The vanishing gradient problem occurs when:
			\begin{enumerate}
				\item Gradients grow too large
				\item Gradients become very small as they propagate through many layers
				\item Weights are too large
				\item Learning rate is too high
			\end{enumerate}
		\end{block}
		\begin{exampleblock}{Answer: B}\end{exampleblock}
		\begin{block}{Explanation}
			Vanishing gradient: gradients shrink exponentially through deep networks with saturating activations (e.g., sigmoid, tanh).
		\end{block}
	\end{frame}
	
	\begin{frame}{S60 --- Exploding Gradient Problem}
		\begin{block}{Question}
			The exploding gradient problem occurs when:
			\begin{enumerate}
				\item Gradients shrink to zero
				\item Gradients grow extremely large across layers
				\item Weights are zero
				\item Learning rate is zero
			\end{enumerate}
		\end{block}
		\begin{exampleblock}{Answer: B}\end{exampleblock}
		\begin{block}{Explanation}
			Exploding gradient: gradients grow exponentially. Remedies include gradient clipping and careful initialization.
		\end{block}
	\end{frame}
	
	\begin{frame}{S61 --- Momentum in Gradient Descent}
		\begin{block}{Question}
			Momentum in gradient descent:
			\begin{enumerate}
				\item Reduces learning rate
				\item Accelerates convergence in consistent directions and dampens oscillation
				\item Eliminates local minima
				\item Converts non-convex to convex
			\end{enumerate}
		\end{block}
		\begin{exampleblock}{Answer: B}\end{exampleblock}
		\begin{block}{Explanation}
			Momentum accumulates velocity, accelerating consistent gradient directions and reducing oscillation.
		\end{block}
	\end{frame}
	
	\begin{frame}{S62 --- Adam Optimizer}
		\begin{block}{Question}
			Adam optimizer combines:
			\begin{enumerate}
				\item Momentum and adaptive learning rates
				\item Batch and stochastic descent
				\item Ridge and LASSO
				\item Precision and recall
			\end{enumerate}
		\end{block}
		\begin{exampleblock}{Answer: A}\end{exampleblock}
		\begin{block}{Explanation}
			Adam = Adaptive Moment Estimation. It combines momentum (first moment) with adaptive per-parameter learning rates (second moment).
		\end{block}
	\end{frame}
	
	\begin{frame}{S63 --- Epoch Definition}
		\begin{block}{Question}
			An epoch in neural network training is:
			\begin{enumerate}
				\item One gradient update
				\item One full pass through the entire training dataset
				\item One validation step
				\item One feature selection step
			\end{enumerate}
		\end{block}
		\begin{exampleblock}{Answer: B}\end{exampleblock}
		\begin{block}{Explanation}
			An epoch = one full pass through the training data. Typically many epochs are required for convergence.
		\end{block}
	\end{frame}
	
	\begin{frame}{S64 --- Batch Gradient Descent}
		\begin{block}{Question}
			Batch gradient descent updates weights using:
			\begin{enumerate}
				\item One random training example
				\item A subset of training examples
				\item The entire training set
				\item No training examples
			\end{enumerate}
		\end{block}
		\begin{exampleblock}{Answer: C}\end{exampleblock}
		\begin{block}{Explanation}
			Batch GD uses the full dataset per update. Stochastic uses one example. Mini-batch uses a subset.
		\end{block}
	\end{frame}
	
	\begin{frame}{S65 --- Stochastic Gradient Descent}
		\begin{block}{Question}
			Stochastic gradient descent (SGD) updates weights using:
			\begin{enumerate}
				\item The full training set
				\item One random training example per step
				\item A fixed subset
				\item No examples
			\end{enumerate}
		\end{block}
		\begin{exampleblock}{Answer: B}\end{exampleblock}
		\begin{block}{Explanation}
			SGD uses one randomly selected example per update. It is noisier but faster per pass.
		\end{block}
	\end{frame}
	
	\begin{frame}{S66 --- Mini-Batch Gradient Descent}
		\begin{block}{Question}
			Mini-batch gradient descent updates weights using:
			\begin{enumerate}
				\item One example
				\item The full dataset
				\item A small subset of training data per step
				\item No data
			\end{enumerate}
		\end{block}
		\begin{exampleblock}{Answer: C}\end{exampleblock}
		\begin{block}{Explanation}
			Mini-batch GD uses a small subset (e.g., 32 or 64 examples). It balances parallelism and stochasticity.
		\end{block}
	\end{frame}
	
	\begin{frame}{S67 --- Overfitting Definition}
		\begin{block}{Question}
			Overfitting occurs when:
			\begin{enumerate}
				\item The model is too simple
				\item The model captures noise in training data and fails to generalize
				\item The model has too few parameters
				\item The model underfits
			\end{enumerate}
		\end{block}
		\begin{exampleblock}{Answer: B}\end{exampleblock}
		\begin{block}{Explanation}
			Overfitting = model captures noise, performs well on training but poorly on test data. Sign: large gap between training and test error.
		\end{block}
	\end{frame}
	
	\begin{frame}{S68 --- Underfitting Definition}
		\begin{block}{Question}
			Underfitting occurs when:
			\begin{enumerate}
				\item The model is too complex
				\item The model fails to capture relevant patterns in the data
				\item The model has too many parameters
				\item The model overfits
			\end{enumerate}
		\end{block}
		\begin{exampleblock}{Answer: B}\end{exampleblock}
		\begin{block}{Explanation}
			Underfitting = model too simple to capture underlying patterns. High bias, low variance.
		\end{block}
	\end{frame}
	
	\begin{frame}{S69 --- Bias-Variance Trade-off}
		\begin{block}{Question}
			The bias-variance trade-off states that:
			\begin{enumerate}
				\item Increasing model complexity always reduces total error
				\item Increasing complexity reduces bias but increases variance
				\item Bias and variance are independent
				\item Variance is always zero
			\end{enumerate}
		\end{block}
		\begin{exampleblock}{Answer: B}\end{exampleblock}
		\begin{block}{Explanation}
			Total error = Bias² + Variance + Irreducible Error. More complex models have lower bias but higher variance.
		\end{block}
	\end{frame}
	
	\begin{frame}{S70 --- Cross-Validation Purpose}
		\begin{block}{Question}
			Cross-validation is used to:
			\begin{enumerate}
				\item Increase training data
				\item Assess model generalization and tune hyperparameters
				\item Remove outliers
				\item Reduce features
			\end{enumerate}
		\end{block}
		\begin{exampleblock}{Answer: B}\end{exampleblock}
		\begin{block}{Explanation}
			Cross-validation splits data into k folds, training on k 1 and validating on the held-out fold, to estimate out-of-sample performance.
		\end{block}
	\end{frame}
	
	\begin{frame}{S71 --- K-Fold Cross-Validation}
		\begin{block}{Question}
			In k-fold cross-validation, the data is:
			\begin{enumerate}
				\item Split into k equal folds; model trained k times, each time holding out a different fold
				\item Split into 2 folds only
				\item Not split
				\item Split randomly without reuse
			\end{enumerate}
		\end{block}
		\begin{exampleblock}{Answer: A}\end{exampleblock}
		\begin{block}{Explanation}
			K-fold CV partitions into k folds; iterate k times, each time training on k-1 folds and validating on the held-out fold.
		\end{block}
	\end{frame}
	
	\begin{frame}{S72 --- Stratified Cross-Validation}
		\begin{block}{Question}
			Stratified cross-validation ensures:
			\begin{enumerate}
				\item Each fold has the same class proportion as the whole dataset
				\item Only the majority class is used
				\item All folds have equal mean
				\item Random assignment is avoided
			\end{enumerate}
		\end{block}
		\begin{exampleblock}{Answer: A}\end{exampleblock}
		\begin{block}{Explanation}
			Stratified CV preserves class ratios in every fold, critical for imbalanced datasets.
		\end{block}
	\end{frame}
	
	\begin{frame}{S73 --- Bootstrapping}
		\begin{block}{Question}
			Bootstrapping in model evaluation is:
			\begin{enumerate}
				\item Sampling with replacement to create new datasets
				\item Sampling without replacement
				\item Splitting data into two halves
				\item Training without validation
			\end{enumerate}
		\end{block}
		\begin{exampleblock}{Answer: A}\end{exampleblock}
		\begin{block}{Explanation}
			Bootstrapping draws samples with replacement. About 37\% of observations are not sampled (out-of-bootstrap).
		\end{block}
	\end{frame}
	
	\begin{frame}{S74 --- Out-of-Bag Fraction}
		\begin{block}{Question}
			In bootstrapping, the expected fraction of observations NOT sampled is approximately:
			\begin{enumerate}
				\item 10\%
				\item 25\%
				\item 37\%
				\item 63\%
			\end{enumerate}
		\end{block}
		\begin{exampleblock}{Answer: C}\end{exampleblock}
		\begin{block}{Explanation}
			Expected out-of-bootstrap fraction = 1 - 1/e = 0.368 = 37\%.
		\end{block}
	\end{frame}
	
	\begin{frame}{S75 --- Grid Search Purpose}
		\begin{block}{Question}
			Grid search is used for:
			\begin{enumerate}
				\item Exhaustively searching through predefined hyperparameter combinations
				\item Removing outliers
				\item Feature selection only
				\item Clustering
			\end{enumerate}
		\end{block}
		\begin{exampleblock}{Answer: A}\end{exampleblock}
		\begin{block}{Explanation}
			Grid search evaluates every combination of hyperparameters within specified ranges, often combined with cross-validation.
		\end{block}
	\end{frame}
	
	\begin{frame}{S76 --- AIC Formula}
		\begin{block}{Question}
			AIC is defined as:
			\begin{enumerate}
				\item $AIC = 2k - 2\ln(L)$
				\item $AIC = k + L$
				\item $AIC = 2k + 2\ln(L)$
				\item $AIC = L/k$
			\end{enumerate}
		\end{block}
		\begin{exampleblock}{Answer: A}\end{exampleblock}
		\begin{block}{Explanation}
			AIC = 2k - 2ln(L), where k = number of parameters, L = likelihood. Lower AIC = better model.
		\end{block}
	\end{frame}
	
	\begin{frame}{S77 --- BIC Formula}
		\begin{block}{Question}
			BIC is defined as:
			\begin{enumerate}
				\item $BIC = k\ln(N) - 2\ln(L)$
				\item $BIC = 2k - 2\ln(L)$
				\item $BIC = k + L$
				\item $BIC = L^2$
			\end{enumerate}
		\end{block}
		\begin{exampleblock}{Answer: A}\end{exampleblock}
		\begin{block}{Explanation}
			BIC = kln(N) - 2ln(L). It penalizes complexity more heavily than AIC for large samples.
		\end{block}
	\end{frame}
	
	\begin{frame}{S78 --- Stepwise Regression}
		\begin{block}{Question}
			Forward stepwise selection:
			\begin{enumerate}
				\item Starts with all features and removes them one by one
				\item Starts with no features and adds them one by one
				\item Uses all features always
				\item Randomly selects features
			\end{enumerate}
		\end{block}
		\begin{exampleblock}{Answer: B}\end{exampleblock}
		\begin{block}{Explanation}
			Forward selection starts with no features and adds the feature that most improves the criterion (e.g., AIC). Backward elimination starts with all features.
		\end{block}
	\end{frame}
	
	\begin{frame}{S79 --- Backward Stepwise Elimination}
		\begin{block}{Question}
			Backward stepwise elimination:
			\begin{enumerate}
				\item Starts with no features and adds them
				\item Starts with all features and removes them one by one
				\item Uses no features
				\item Randomly selects features
			\end{enumerate}
		\end{block}
		\begin{exampleblock}{Answer: B}\end{exampleblock}
		\begin{block}{Explanation}
			Backward elimination starts with the full model and removes the feature whose removal most improves the criterion.
		\end{block}
	\end{frame}
	
	% =================================================================
	\section{Natural Language Processing (Additional Topics)}
	% =================================================================
	
	\begin{frame}{S80 --- Tokenization Definition}
		\begin{block}{Question}
			Tokenization in NLP is:
			\begin{enumerate}
				\item Removing stop words
				\item Separating text into sentences and words
				\item Converting words to vectors
				\item Replacing words with synonyms
			\end{enumerate}
		\end{block}
		\begin{exampleblock}{Answer: B}\end{exampleblock}
		\begin{block}{Explanation}
			Tokenization separates documents into sentences (sentence tokenization) and sentences into words (word tokenization).
		\end{block}
	\end{frame}
	
	\begin{frame}{S81 --- Stop Word Removal}
		\begin{block}{Question}
			Stop word removal involves:
			\begin{enumerate}
				\item Removing words with no informational value (e.g., ``the'', ``a'', ``and'')
				\item Converting words to vectors
				\item Separating sentences
				\item Replacing words with roots
			\end{enumerate}
		\end{block}
		\begin{exampleblock}{Answer: A}\end{exampleblock}
		\begin{block}{Explanation}
			Stop words are common words with little informational value removed to improve analysis efficiency.
		\end{block}
	\end{frame}
	
	\begin{frame}{S82 --- Stemming Definition}
		\begin{block}{Question}
			Stemming in NLP:
			\begin{enumerate}
				\item Replaces words with their dictionary form
				\item Replaces words with their root form by cutting prefixes and suffixes
				\item Removes stop words
				\item Converts to vectors
			\end{enumerate}
		\end{block}
		\begin{exampleblock}{Answer: B}\end{exampleblock}
		\begin{block}{Explanation}
			Stemming cuts prefixes/suffixes (e.g., disappointing → disappoint). Lemmatization uses dictionary forms.
		\end{block}
	\end{frame}
	
	\begin{frame}{S83 --- Lemmatization Definition}
		\begin{block}{Question}
			Lemmatization in NLP:
			\begin{enumerate}
				\item Cuts prefixes and suffixes
				\item Replaces words with their dictionary form based on meaning and part of speech
				\item Removes stop words
				\item Converts to vectors
			\end{enumerate}
		\end{block}
		\begin{exampleblock}{Answer: B}\end{exampleblock}
		\begin{block}{Explanation}
			Lemmatization uses vocabulary and morphological analysis to find the lemma (dictionary form) of a word.
		\end{block}
	\end{frame}
	
	\begin{frame}{S84 --- Bag of Words (Binary)}
		\begin{block}{Question}
			In binary bag of words, each element of the document vector is:
			\begin{enumerate}
				\item The count of word occurrences
				\item 1 if the word appears, 0 otherwise
				\item The TF-IDF score
				\item The word embedding
			\end{enumerate}
		\end{block}
		\begin{exampleblock}{Answer: B}\end{exampleblock}
		\begin{block}{Explanation}
			Binary BoW: 1 = word present, 0 = absent. Count BoW records the number of occurrences.
		\end{block}
	\end{frame}
	
	\begin{frame}{S85 --- Bag of Words (Count)}
		\begin{block}{Question}
			In count bag of words, each element of the document vector is:
			\begin{enumerate}
				\item 1 or 0
				\item The number of times the word appears in the document
				\item The TF-IDF score
				\item The word embedding
			\end{enumerate}
		\end{block}
		\begin{exampleblock}{Answer: B}\end{exampleblock}
		\begin{block}{Explanation}
			Count BoW records the frequency of each word in the document, producing non-negative integers.
		\end{block}
	\end{frame}
	
	\begin{frame}{S86 --- TF-IDF Formula}
		\begin{block}{Question}
			The TF-IDF weight of term i in document j is:
			\begin{enumerate}
				\item $TF_{ij} + IDF_i$
				\item $TF_{ij} \times IDF_i$
				\item $TF_{ij} / IDF_i$
				\item $TF_{ij} - IDF_i$
			\end{enumerate}
		\end{block}
		\begin{exampleblock}{Answer: B}\end{exampleblock}
		\begin{block}{Explanation}
			TF-IDF = term frequency × inverse document frequency. It weights words that are frequent in a document but rare across the corpus.
		\end{block}
	\end{frame}
	
	\begin{frame}{S87 --- IDF Formula}
		\begin{block}{Question}
			The inverse document frequency of term i across D documents is:
			\begin{enumerate}
				\item $\ln(D/d_i)$
				\item $D \cdot d_i$
				\item $D + d_i$
				\item $D - d_i$
			\end{enumerate}
		\end{block}
		\begin{exampleblock}{Answer: A}\end{exampleblock}
		\begin{block}{Explanation}
			IDFᵢ = ln(D/dᵢ), where dᵢ = number of documents containing term i. If term appears in all documents, IDF = 0.
		\end{block}
	\end{frame}
	
	\begin{frame}{S88 --- N-gram Definition}
		\begin{block}{Question}
			An n-gram in NLP is:
			\begin{enumerate}
				\item A single word
				\item n contiguous words treated as a single unit
				\item A sentence
				\item A document
			\end{enumerate}
		\end{block}
		\begin{exampleblock}{Answer: B}\end{exampleblock}
		\begin{block}{Explanation}
			An n-gram is n contiguous words treated as a unit. Uni-gram = 1 word, bi-gram = 2 words, etc.
		\end{block}
	\end{frame}
	
	\begin{frame}{S89 --- L2 Vector Normalization}
		\begin{block}{Question}
			L2 normalization of a vector v divides each element by:
			\begin{enumerate}
				\item $\sum v_i$
				\item $\sqrt{\sum v_i^2}$
				\item $\max(v_i)$
				\item $|v_1|$
			\end{enumerate}
		\end{block}
		\begin{exampleblock}{Answer: B}\end{exampleblock}
		\begin{block}{Explanation}
			L2 norm = √(Σvᵢ²). Dividing by it produces a unit vector.
		\end{block}
	\end{frame}
	
	\begin{frame}{S90 --- Naive Bayes Assumption}
		\begin{block}{Question}
			The ``naive'' in Naive Bayes refers to:
			\begin{enumerate}
				\item Assuming all words are independent
				\item Assuming all words are dependent
				\item Assuming equal class priors
				\item Assuming no training data
			\end{enumerate}
		\end{block}
		\begin{exampleblock}{Answer: A}\end{exampleblock}
		\begin{block}{Explanation}
			Naive Bayes assumes features (words) are independent of each other, which simplifies computation but is rarely true.
		\end{block}
	\end{frame}
	
	\begin{frame}{S91 --- Naive Bayes Classification Rule}
		\begin{block}{Question}
			The Naive Bayes classification rule is:
			\begin{enumerate}
				\item $\hat{c} = \arg\max_c P(c)$
				\item $\hat{c} = \arg\max_c P(w_1|c)P(w_2|c)\cdots P(w_N|c)P(c)$
				\item $\hat{c} = \arg\min_c P(c)$
				\item $\hat{c} = \sum_c P(c)$
			\end{enumerate}
		\end{block}
		\begin{exampleblock}{Answer: B}\end{exampleblock}
		\begin{block}{Explanation}
			Naive Bayes predicts the class maximizing the product of conditional word probabilities and class prior.
		\end{block}
	\end{frame}
	
	\begin{frame}{S92 --- Laplace Smoothing Purpose}
		\begin{block}{Question}
			Laplace smoothing in Naive Bayes addresses:
			\begin{enumerate}
				\item Computational overflow
				\item Zero probability problem
				\item Unknown words
				\item Class imbalance
			\end{enumerate}
		\end{block}
		\begin{exampleblock}{Answer: B}\end{exampleblock}
		\begin{block}{Explanation}
			Laplace smoothing adds a small positive integer (λ=1) to all counts, preventing zero probabilities when a word never appears in a class.
		\end{block}
	\end{frame}
	
	\begin{frame}{S93 --- Bayes Rule}
		\begin{block}{Question}
			Bayes' rule is:
			\begin{enumerate}
				\item $P(c|d) = P(d|c)P(c)/P(d)$
				\item $P(c|d) = P(c) + P(d)$
				\item $P(c|d) = P(d|c)/P(c)$
				\item $P(c|d) = P(c) \cdot P(d)$
			\end{enumerate}
		\end{block}
		\begin{exampleblock}{Answer: A}\end{exampleblock}
		\begin{block}{Explanation}
			Bayes: posterior = likelihood × prior / evidence. Used in Naive Bayes classification.
		\end{block}
	\end{frame}
	
	\begin{frame}{S94 --- Word2Vec Architectures}
		\begin{block}{Question}
			The two architectures of Word2Vec are:
			\begin{enumerate}
				\item CBoW and Skip-gram
				\item RNN and LSTM
				\item BERT and GPT
				\item CNN and RNN
			\end{enumerate}
		\end{block}
		\begin{exampleblock}{Answer: A}\end{exampleblock}
		\begin{block}{Explanation}
			Word2Vec has two architectures: Continuous Bag of Words (CBoW) and Skip-gram.
		\end{block}
	\end{frame}
	
	\begin{frame}{S95 --- CBoW vs Skip-gram}
		\begin{block}{Question}
			CBoW predicts   from  , while Skip-gram predicts   from  .
			\begin{enumerate}
				\item Context words; target word; target word; context words
				\item Target word; context words; context words; target word
				\item Sentences; words; words; sentences
				\item Documents; words; words; documents
			\end{enumerate}
		\end{block}
		\begin{exampleblock}{Answer: B}\end{exampleblock}
		\begin{block}{Explanation}
			CBoW: predict target word from context words. Skip-gram: predict context words from target word.
		\end{block}
	\end{frame}
	
	\begin{frame}{S96 --- Cosine Similarity Formula}
		\begin{block}{Question}
			Cosine similarity between vectors P and Q is:
			\begin{enumerate}
				\item $\|P - Q\|$
				\item $P \cdot Q / (\|P\| \|Q\|)$
				\item $\|P\| + \|Q\|$
				\item $P/Q$
			\end{enumerate}
		\end{block}
		\begin{exampleblock}{Answer: B}\end{exampleblock}
		\begin{block}{Explanation}
			Cosine similarity = dot product / product of magnitudes. Range is [-1, 1]. Used in Word2Vec for word similarity.
		\end{block}
	\end{frame}
	
	\begin{frame}{S97 --- GloVe Full Name}
		\begin{block}{Question}
			GloVe stands for:
			\begin{enumerate}
				\item Global Vectors for Word Representation
				\item Global Learning of Vector Embeddings
				\item Generalized Local Vector Encoding
				\item Gradient Loss Vector Estimation
			\end{enumerate}
		\end{block}
		\begin{exampleblock}{Answer: A}\end{exampleblock}
		\begin{block}{Explanation}
			GloVe = Global Vectors for Word Representation, developed at Stanford. It is an alternative to Word2Vec.
		\end{block}
	\end{frame}
	
	\begin{frame}{S98 --- RNN Hidden State}
		\begin{block}{Question}
			In an RNN, the hidden state at time t depends on:
			\begin{enumerate}
				\item Only the current input
				\item Only the previous hidden state
				\item Both the current input and the previous hidden state
				\item Neither
			\end{enumerate}
		\end{block}
		\begin{exampleblock}{Answer: C}\end{exampleblock}
		\begin{block}{Explanation}
			RNN hidden state combines current input and previous hidden state, enabling memory of past information.
		\end{block}
	\end{frame}
	
	\begin{frame}{S99 --- LSTM Gates}
		\begin{block}{Question}
			An LSTM cell typically contains how many gates?
			\begin{enumerate}
				\item 1
				\item 2
				\item 3
				\item 4
			\end{enumerate}
		\end{block}
		\begin{exampleblock}{Answer: C}\end{exampleblock}
		\begin{block}{Explanation}
			LSTM has three gates: forget gate (discard), input gate (store), output gate (output). They regulate information flow.
		\end{block}
	\end{frame}
	
	\begin{frame}{S100 --- LSTM Forget Gate}
		\begin{block}{Question}
			The forget gate in an LSTM determines:
			\begin{enumerate}
				\item Which part of memory is discarded
				\item Which part of output is stored
				\item Which part of memory is output
				\item The learning rate
			\end{enumerate}
		\end{block}
		\begin{exampleblock}{Answer: A}\end{exampleblock}
		\begin{block}{Explanation}
			Forget gate decides what to discard from memory. Input gate decides what to store. Output gate decides what to output.
		\end{block}
	\end{frame}
	
	\begin{frame}{S101 --- Transformer Paper}
		\begin{block}{Question}
			The 2017 paper introducing the transformer architecture is titled:
			\begin{enumerate}
				\item ``Attention Is All You Need''
				\item ``Deep Learning for AI''
				\item ``Neural Networks and Learning''
				\item ``Machine Learning Fundamentals''
			\end{enumerate}
		\end{block}
		\begin{exampleblock}{Answer: A}\end{exampleblock}
		\begin{block}{Explanation}
			``Attention Is All You Need'' by Vaswani et al. (2017) introduced the transformer based on self-attention.
		\end{block}
	\end{frame}
	
	\begin{frame}{S102 --- BERT Full Name}
		\begin{block}{Question}
			BERT stands for:
			\begin{enumerate}
				\item Bidirectional Encoder Representations from Transformers
				\item Bayesian Encoder for Recursive Tasks
				\item Binary Encoder Representation Technique
				\item Bidirectional Encoder for Random Text
			\end{enumerate}
		\end{block}
		\begin{exampleblock}{Answer: A}\end{exampleblock}
		\begin{block}{Explanation}
			BERT = Bidirectional Encoder Representations from Transformers. Encoder-only model released by Google in 2018.
		\end{block}
	\end{frame}
	
	\begin{frame}{S103 --- GPT Full Name}
		\begin{block}{Question}
			GPT stands for:
			\begin{enumerate}
				\item General Purpose Transformer
				\item Generative Pre-trained Transformer
				\item Gradient Pre-training Technique
				\item Global Parameter Training
			\end{enumerate}
		\end{block}
		\begin{exampleblock}{Answer: B}\end{exampleblock}
		\begin{block}{Explanation}
			GPT = Generative Pre-trained Transformer. Decoder-only autoregressive model from OpenAI.
		\end{block}
	\end{frame}
	
	\begin{frame}{S104 --- BART Full Name}
		\begin{block}{Question}
			BART stands for:
			\begin{enumerate}
				\item Bidirectional and Auto-Regressive Transformer
				\item Bayesian Auto-Regressive Technique
				\item Binary Auto-Regressive Transformer
				\item Bidirectional Attention for Recognition Tasks
			\end{enumerate}
		\end{block}
		\begin{exampleblock}{Answer: A}\end{exampleblock}
		\begin{block}{Explanation}
			BART = Bidirectional and Auto-Regressive Transformer. Combines encoder and decoder for understanding and generation.
		\end{block}
	\end{frame}
	
	\begin{frame}{S105 --- LSTM Full Name}
		\begin{block}{Question}
			LSTM stands for:
			\begin{enumerate}
				\item Long Short-Term Memory
				\item Linear Sequential Time Model
				\item Layered Stochastic Transformation Model
				\item Low-Signal Temporal Method
			\end{enumerate}
		\end{block}
		\begin{exampleblock}{Answer: A}\end{exampleblock}
		\begin{block}{Explanation}
			LSTM = Long Short-Term Memory. Type of RNN that mitigates vanishing gradient via gating.
		\end{block}
	\end{frame}
	
	\begin{frame}{S106 --- Attention Score Formula}
		\begin{block}{Question}
			The scaled dot-product attention score between query Qᵢ and key Kⱼ is:
			\begin{enumerate}
				\item $Q_i \cdot K_j$
				\item $Q_i \cdot K_j / \sqrt{d_k}$
				\item $(Q_i + K_j)/2$
				\item $|Q_i - K_j|$
			\end{enumerate}
		\end{block}
		\begin{exampleblock}{Answer: B}\end{exampleblock}
		\begin{block}{Explanation}
			Scaled dot-product attention divides by √dₖ (key dimension) to prevent softmax saturation.
		\end{block}
	\end{frame}
	
	\begin{frame}{S107 --- Softmax Function}
		\begin{block}{Question}
			The softmax function is:
			\begin{enumerate}
				\item $\text{softmax}(x_i) = e^{x_i}/\sum_j e^{x_j}$
				\item $\text{softmax}(x_i) = x_i/\sum_j x_j$
				\item $\text{softmax}(x_i) = e^{x_i}$
				\item $\text{softmax}(x_i) = \ln(x_i)$
			\end{enumerate}
		\end{block}
		\begin{exampleblock}{Answer: A}\end{exampleblock}
		\begin{block}{Explanation}
			Softmax normalizes exponentials into a probability distribution. Used in attention and multiclass outputs.
		\end{block}
	\end{frame}
	
	\begin{frame}{S108 --- Context Length Definition}
		\begin{block}{Question}
			Context length (context window) of an LLM is:
			\begin{enumerate}
				\item The number of parameters
				\item The maximum number of tokens the model can process in a single input/output
				\item The training time
				\item The vocabulary size
			\end{enumerate}
		\end{block}
		\begin{exampleblock}{Answer: B}\end{exampleblock}
		\begin{block}{Explanation}
			Context length = maximum tokens processable in one sequence. GPT-4 has 32k; Gemini 2.5 Pro has 1048k.
		\end{block}
	\end{frame}
	
	\begin{frame}{S109 --- Statelessness of LLMs}
		\begin{block}{Question}
			LLMs are stateless, meaning:
			\begin{enumerate}
				\item They retain calculations between prompts
				\item They do not retain calculations or information between processing one prompt and the next
				\item They store all conversations
				\item They learn from every interaction
			\end{enumerate}
		\end{block}
		\begin{exampleblock}{Answer: B}\end{exampleblock}
		\begin{block}{Explanation}
			Statelessness = no memory between prompts. Users must feed previous context if needed.
		\end{block}
	\end{frame}
	
	\begin{frame}{S110 --- Temperature Parameter}
		\begin{block}{Question}
			Temperature in LLMs controls:
			\begin{enumerate}
				\item The number of tokens generated
				\item The shape of the probability distribution (randomness/creativity)
				\item The context length
				\item The learning rate
			\end{enumerate}
		\end{block}
		\begin{exampleblock}{Answer: B}\end{exampleblock}
		\begin{block}{Explanation}
			Temperature < 1 narrows distribution (more predictable). Temperature > 1 flattens distribution (more creative/random).
		\end{block}
	\end{frame}
	
	\begin{frame}{S111 --- Top-K Sampling}
		\begin{block}{Question}
			Top-K sampling restricts the model to consider:
			\begin{enumerate}
				\item All tokens
				\item The K most probable tokens at each step
				\item The K least probable tokens
				\item Random tokens
			\end{enumerate}
		\end{block}
		\begin{exampleblock}{Answer: B}\end{exampleblock}
		\begin{block}{Explanation}
			Top-K sampling considers only the K most probable tokens. Low K sacrifices creativity; high K allows more variety.
		\end{block}
	\end{frame}
	
	\begin{frame}{S112 --- Top-P Sampling}
		\begin{block}{Question}
			Top-P (nucleus) sampling selects:
			\begin{enumerate}
				\item The K most probable tokens
				\item The smallest set of tokens whose cumulative probability ≥ P
				\item All tokens
				\item Random tokens
			\end{enumerate}
		\end{block}
		\begin{exampleblock}{Answer: B}\end{exampleblock}
		\begin{block}{Explanation}
			Top-P selects the smallest set of tokens with cumulative probability ≥ P. A low P sacrifices creativity.
		\end{block}
	\end{frame}
	
	\begin{frame}{S113 --- Chain-of-Thought Prompting}
		\begin{block}{Question}
			Chain-of-Thought (CoT) prompting:
			\begin{enumerate}
				\item Asks for direct answers
				\item Prompts with intermediate reasoning steps
				\item Uses random prompts
				\item Removes context
			\end{enumerate}
		\end{block}
		\begin{exampleblock}{Answer: B}\end{exampleblock}
		\begin{block}{Explanation}
			CoT prompting encourages step-by-step reasoning, improving performance on multi-step tasks like arithmetic.
		\end{block}
	\end{frame}
	
	\begin{frame}{S114 --- Zero-Shot Prompting}
		\begin{block}{Question}
			Zero-shot prompting involves:
			\begin{enumerate}
				\item Providing one example
				\item Providing multiple examples
				\item Providing no examples, just the question
				\item Providing only context
			\end{enumerate}
		\end{block}
		\begin{exampleblock}{Answer: C}\end{exampleblock}
		\begin{block}{Explanation}
			Zero-shot = no examples. One-shot = one example. Few-shot = multiple examples.
		\end{block}
	\end{frame}
	
	\begin{frame}{S115 --- Few-Shot Prompting}
		\begin{block}{Question}
			Few-shot prompting involves:
			\begin{enumerate}
				\item No examples
				\item One example
				\item Multiple examples
				\item Only context
			\end{enumerate}
		\end{block}
		\begin{exampleblock}{Answer: C}\end{exampleblock}
		\begin{block}{Explanation}
			Few-shot = multiple relevant examples provided to guide the model's response.
		\end{block}
	\end{frame}
	
	\begin{frame}{S116 --- RAG Full Name}
		\begin{block}{Question}
			RAG stands for:
			\begin{enumerate}
				\item Retrieval Augmented Generation
				\item Random Attention Gate
				\item Recurrent Adversarial Graph
				\item Regularized Attention Gradient
			\end{enumerate}
		\end{block}
		\begin{exampleblock}{Answer: A}\end{exampleblock}
		\begin{block}{Explanation}
			RAG = Retrieval Augmented Generation. Augments LLM responses with retrieved documents for factual accuracy.
		\end{block}
	\end{frame}
	
	\begin{frame}{S117 --- RLHF Full Name}
		\begin{block}{Question}
			RLHF stands for:
			\begin{enumerate}
				\item Reinforcement Learning from Human Feedback
				\item Recursive Layered Hidden Features
				\item Random Loss with High Frequency
				\item Regularized Linear Hierarchical Framework
			\end{enumerate}
		\end{block}
		\begin{exampleblock}{Answer: A}\end{exampleblock}
		\begin{block}{Explanation}
			RLHF = Reinforcement Learning from Human Feedback. Used in LLM training to align outputs with human preferences.
		\end{block}
	\end{frame}
	
	\begin{frame}{S118 --- Hallucination Definition}
		\begin{block}{Question}
			Hallucination in GenAI refers to:
			\begin{enumerate}
				\item Generating factually incorrect but plausible output
				\item Generating random noise
				\item Generating too many tokens
				\item Generating too few tokens
			\end{enumerate}
		\end{block}
		\begin{exampleblock}{Answer: A}\end{exampleblock}
		\begin{block}{Explanation}
			Hallucination = factually incorrect, misleading, or fabricated output that appears credible.
		\end{block}
	\end{frame}
	
	\begin{frame}{S119 --- MMLU Benchmark}
		\begin{block}{Question}
			MMLU stands for:
			\begin{enumerate}
				\item Massive Multitask Language Understanding
				\item Multi-Modal Language Understanding
				\item Machine Learning Utility
				\item Multi-Layer Language Understanding
			\end{enumerate}
		\end{block}
		\begin{exampleblock}{Answer: A}\end{exampleblock}
		\begin{block}{Explanation}
			MMLU = Massive Multitask Language Understanding. Evaluates models across 57 subjects with multiple-choice questions.
		\end{block}
	\end{frame}
	
	\begin{frame}{S120 --- HELM Full Name}
		\begin{block}{Question}
			HELM stands for:
			\begin{enumerate}
				\item Holistic Evaluation of Language Models
				\item Hierarchical Evaluation of Learning Models
				\item Hybrid Ensemble Learning Method
				\item High-Efficiency Language Model
			\end{enumerate}
		\end{block}
		\begin{exampleblock}{Answer: A}\end{exampleblock}
		\begin{block}{Explanation}
			HELM = Holistic Evaluation of Language Models, a leaderboard from Stanford CRFM.
		\end{block}
	\end{frame}
	
	% =================================================================
	\section{Reinforcement Learning (Additional Topics)}
	% =================================================================
	
	\begin{frame}{S121 --- Agent Definition}
		\begin{block}{Question}
			In reinforcement learning, the agent is:
			\begin{enumerate}
				\item The environment
				\item The learner or decision-maker
				\item The reward
				\item The state
			\end{enumerate}
		\end{block}
		\begin{exampleblock}{Answer: B}\end{exampleblock}
		\begin{block}{Explanation}
			Agent = learner/decision-maker. In MAB, the agent is the gambler.
		\end{block}
	\end{frame}
	
	\begin{frame}{S122 --- Reward Definition}
		\begin{block}{Question}
			In reinforcement learning, reward is:
			\begin{enumerate}
				\item The state
				\item The action
				\item Feedback from the environment after taking an action
				\item The policy
			\end{enumerate}
		\end{block}
		\begin{exampleblock}{Answer: C}\end{exampleblock}
		\begin{block}{Explanation}
			Reward = feedback (positive or negative) from the environment after an action. The agent aims to maximize cumulative reward.
		\end{block}
	\end{frame}
	
	\begin{frame}{S123 --- Policy Definition}
		\begin{block}{Question}
			In reinforcement learning, a policy is:
			\begin{enumerate}
				\item The reward function
				\item The agent's strategy for choosing actions given states
				\item The state space
				\item The action space
			\end{enumerate}
		\end{block}
		\begin{exampleblock}{Answer: B}\end{exampleblock}
		\begin{block}{Explanation}
			Policy (π) = strategy mapping states to actions. Goal: find optimal policy maximizing long-term reward.
		\end{block}
	\end{frame}
	
	\begin{frame}{S124 --- Value Function Definition}
		\begin{block}{Question}
			The value function V(s) in reinforcement learning estimates:
			\begin{enumerate}
				\item Immediate reward
				\item Total expected future reward from state s following a policy
				\item The number of states
				\item The action space size
			\end{enumerate}
		\end{block}
		\begin{exampleblock}{Answer: B}\end{exampleblock}
		\begin{block}{Explanation}
			V(s) = expected cumulative future reward starting from state s and following a policy.
		\end{block}
	\end{frame}
	
	\begin{frame}{S125 --- Action-Value Function Definition}
		\begin{block}{Question}
			The action-value function Q(s,a) estimates:
			\begin{enumerate}
				\item Value of state s only
				\item Value of taking action a in state s
				\item Number of actions
				\item Reward only
			\end{enumerate}
		\end{block}
		\begin{exampleblock}{Answer: B}\end{exampleblock}
		\begin{block}{Explanation}
			Q(s,a) = expected cumulative future reward from taking action a in state s, then following a policy.
		\end{block}
	\end{frame}
	
	\begin{frame}{S126 --- Epsilon-Greedy Strategy}
		\begin{block}{Question}
			The ε-greedy strategy balances:
			\begin{enumerate}
				\item Only exploration
				\item Only exploitation
				\item Exploration (random action with probability ε) and exploitation (best action with probability 1-ε)
				\item Neither
			\end{enumerate}
		\end{block}
		\begin{exampleblock}{Answer: C}\end{exampleblock}
		\begin{block}{Explanation}
			ε-greedy: with probability ε explore randomly; with probability 1-ε exploit the best-known action.
		\end{block}
	\end{frame}
	
	\begin{frame}{S127 --- Greedy Strategy}
		\begin{block}{Question}
			The greedy strategy in MAB:
			\begin{enumerate}
				\item Always explores
				\item Always chooses the action with the best reward so far
				\item Chooses randomly
				\item Never chooses
			\end{enumerate}
		\end{block}
		\begin{exampleblock}{Answer: B}\end{exampleblock}
		\begin{block}{Explanation}
			Greedy = pure exploitation. It always chooses the action with the highest estimated value, ignoring exploration.
		\end{block}
	\end{frame}
	
	\begin{frame}{S128 --- Random Strategy}
		\begin{block}{Question}
			The random strategy in MAB:
			\begin{enumerate}
				\item Exploits best action
				\item Randomly selects an action
				\item Never selects
				\item Uses gradient descent
			\end{enumerate}
		\end{block}
		\begin{exampleblock}{Answer: B}\end{exampleblock}
		\begin{block}{Explanation}
			Random strategy = pure exploration. It randomly selects actions without using past knowledge.
		\end{block}
	\end{frame}
	
	\begin{frame}{S129 --- Monte Carlo Method}
		\begin{block}{Question}
			The Monte Carlo method in RL updates value estimates:
			\begin{enumerate}
				\item After each step
				\item Only at the end of a complete episode
				\item Never
				\item Randomly
			\end{enumerate}
		\end{block}
		\begin{exampleblock}{Answer: B}\end{exampleblock}
		\begin{block}{Explanation}
			Monte Carlo updates only at episode end using actual returns. High variance, unbiased.
		\end{block}
	\end{frame}
	
	\begin{frame}{S130 --- Temporal Difference Method}
		\begin{block}{Question}
			The Temporal Difference (TD) method in RL updates value estimates:
			\begin{enumerate}
				\item Only at episode end
				\item After each step using bootstrapping
				\item Never
				\item Only at the beginning
			\end{enumerate}
		\end{block}
		\begin{exampleblock}{Answer: B}\end{exampleblock}
		\begin{block}{Explanation}
			TD updates after each step by bootstrapping: using current estimate of next state's value. Lower variance than MC.
		\end{block}
	\end{frame}
	
	\begin{frame}{S131 --- Q-Learning}
		\begin{block}{Question}
			Q-learning is a:
			\begin{enumerate}
				\item Monte Carlo method
				\item Temporal Difference method
				\item Supervised learning method
				\item Unsupervised learning method
			\end{enumerate}
		\end{block}
		\begin{exampleblock}{Answer: B}\end{exampleblock}
		\begin{block}{Explanation}
			Q-learning is a TD method that learns action-value function Q(s,a) step-by-step.
		\end{block}
	\end{frame}
	
	\begin{frame}{S132 --- DQN Full Name}
		\begin{block}{Question}
			DQN stands for:
			\begin{enumerate}
				\item Deep Q-Network
				\item Dynamic Quantum Network
				\item Dual Q-Network
				\item Distributed Q-Network
			\end{enumerate}
		\end{block}
		\begin{exampleblock}{Answer: A}\end{exampleblock}
		\begin{block}{Explanation}
			DQN = Deep Q-Network. Combines TD method with neural networks to approximate Q(s,a) for large state spaces.
		\end{block}
	\end{frame}
	
	\begin{frame}{S133 --- DQN Key Components}
		\begin{block}{Question}
			DQN consists of which three key components?
			\begin{enumerate}
				\item Policy network, experience replay, target network
				\item Encoder, decoder, attention
				\item Input, hidden, output
				\item Actor, critic, reward
			\end{enumerate}
		\end{block}
		\begin{exampleblock}{Answer: A}\end{exampleblock}
		\begin{block}{Explanation}
			DQN: policy network (Q-network), experience replay (memory buffer), target network (stable Q-values).
		\end{block}
	\end{frame}
	
	\begin{frame}{S134 --- Experience Replay Purpose}
		\begin{block}{Question}
			Experience replay in DQN:
			\begin{enumerate}
				\item Increases correlation between samples
				\item Breaks temporal correlation by random sampling of past experiences
				\item Removes outliers
				\item Increases learning rate
			\end{enumerate}
		\end{block}
		\begin{exampleblock}{Answer: B}\end{exampleblock}
		\begin{block}{Explanation}
			Experience replay stores past experiences and samples mini-batches randomly, breaking temporal correlation.
		\end{block}
	\end{frame}
	
	\begin{frame}{S135 --- Policy-Based vs Value-Based RL}
		\begin{block}{Question}
			Policy-based RL methods:
			\begin{enumerate}
				\item Learn value functions first
				\item Learn the policy directly
				\item Learn neither
				\item Only work with discrete actions
			\end{enumerate}
		\end{block}
		\begin{exampleblock}{Answer: B}\end{exampleblock}
		\begin{block}{Explanation}
			Policy-based methods learn the policy directly (parameterized by θ), suitable for continuous action spaces.
		\end{block}
	\end{frame}
	
	\begin{frame}{S136 --- Actor-Critic Methods}
		\begin{block}{Question}
			Actor-Critic methods combine:
			\begin{enumerate}
				\item Value-based and policy-based approaches
				\item Supervised and unsupervised learning
				\item CNN and RNN
				\item Bagging and boosting
			\end{enumerate}
		\end{block}
		\begin{exampleblock}{Answer: A}\end{exampleblock}
		\begin{block}{Explanation}
			Actor-Critic: actor (policy) + critic (value function). Combines strengths of policy-based and value-based methods.
		\end{block}
	\end{frame}
	
	% =================================================================
	\section{Ethics and Governance (Additional Topics)}
	% =================================================================
	
	\begin{frame}{S137 --- Utilitarianism Founders}
		\begin{block}{Question}
			The founders of utilitarianism are:
			\begin{enumerate}
				\item Kant and Aristotle
				\item Bentham and Mill
				\item Rawls and Nozick
				\item Plato and Socrates
			\end{enumerate}
		\end{block}
		\begin{exampleblock}{Answer: B}\end{exampleblock}
		\begin{block}{Explanation}
			Jeremy Bentham and John Stuart Mill developed classical utilitarianism in the 18th–19th centuries.
		\end{block}
	\end{frame}
	
	\begin{frame}{S138 --- Categorical Imperative}
		\begin{block}{Question}
			The categorical imperative is attributed to:
			\begin{enumerate}
				\item Immanuel Kant
				\item John Stuart Mill
				\item Jeremy Bentham
				\item Aristotle
			\end{enumerate}
		\end{block}
		\begin{exampleblock}{Answer: A}\end{exampleblock}
		\begin{block}{Explanation}
			Kant's categorical imperative: act only on maxims that could be universal laws; treat people as ends, not means.
		\end{block}
	\end{frame}
	
	\begin{frame}{S139 --- Virtue Ethics Origin}
		\begin{block}{Question}
			Virtue ethics is most closely associated with:
			\begin{enumerate}
				\item Aristotle
				\item Kant
				\item Mill
				\item Bentham
			\end{enumerate}
		\end{block}
		\begin{exampleblock}{Answer: A}\end{exampleblock}
		\begin{block}{Explanation}
			Virtue ethics has roots in Aristotle's Nicomachean Ethics, emphasizing virtuous character traits.
		\end{block}
	\end{frame}
	
	\begin{frame}{S140 --- Nonmaleficence Definition}
		\begin{block}{Question}
			Nonmaleficence refers to:
			\begin{enumerate}
				\item Actively promoting welfare
				\item Avoiding harm and unjustified risk
				\item Fairness
				\item Respecting autonomy
			\end{enumerate}
		\end{block}
		\begin{exampleblock}{Answer: B}\end{exampleblock}
		\begin{block}{Explanation}
			Nonmaleficence = do no harm. It forbids imposing unjustified risks.
		\end{block}
	\end{frame}
	
	\begin{frame}{S141 --- Beneficence Definition}
		\begin{block}{Question}
			Beneficence refers to:
			\begin{enumerate}
				\item Avoiding harm
				\item Actively promoting the welfare of others
				\item Respecting autonomy
				\item Ensuring fairness
			\end{enumerate}
		\end{block}
		\begin{exampleblock}{Answer: B}\end{exampleblock}
		\begin{block}{Explanation}
			Beneficence = do good. It requires positive action to help others.
		\end{block}
	\end{frame}
	
	\begin{frame}{S142 --- Autonomy Definition}
		\begin{block}{Question}
			Autonomy as an AI ethics principle requires:
			\begin{enumerate}
				\item Respecting individual decision-making
				\item Ensuring fairness
				\item Avoiding harm
				\item Benefiting others
			\end{enumerate}
		\end{block}
		\begin{exampleblock}{Answer: A}\end{exampleblock}
		\begin{block}{Explanation}
			Autonomy = respect for individual self-determination and consent.
		\end{block}
	\end{frame}
	
	\begin{frame}{S143 --- Justice Definition}
		\begin{block}{Question}
			Justice as an AI ethics principle primarily concerns:
			\begin{enumerate}
				\item Fairness, equality, and impartiality
				\item Avoiding harm
				\item Benefiting others
				\item Respecting autonomy
			\end{enumerate}
		\end{block}
		\begin{exampleblock}{Answer: A}\end{exampleblock}
		\begin{block}{Explanation}
			Justice = fairness, equality, impartiality. Distributive justice concerns allocation; procedural justice concerns fair process.
		\end{block}
	\end{frame}
	
	\begin{frame}{S144 --- Consequentialism Definition}
		\begin{block}{Question}
			Consequentialism judges actions by:
			\begin{enumerate}
				\item Adherence to rules
				\item Outcomes or consequences
				\item Character traits
				\item Legal compliance
			\end{enumerate}
		\end{block}
		\begin{exampleblock}{Answer: B}\end{exampleblock}
		\begin{block}{Explanation}
			Consequentialism judges by outcomes. Utilitarianism is the most common form, maximizing overall welfare.
		\end{block}
	\end{frame}
	
	\begin{frame}{S145 --- Deontology Definition}
		\begin{block}{Question}
			Deontology judges actions by:
			\begin{enumerate}
				\item Outcomes
				\item Adherence to duties and rules
				\item Character
				\item Profit
			\end{enumerate}
		\end{block}
		\begin{exampleblock}{Answer: B}\end{exampleblock}
		\begin{block}{Explanation}
			Deontology judges by duties and rules irrespective of consequences. Kant is the primary proponent.
		\end{block}
	\end{frame}
	
	\begin{frame}{S146 --- Virtue Ethics Definition}
		\begin{block}{Question}
			Virtue ethics focuses on:
			\begin{enumerate}
				\item Outcomes
				\item Rules
				\item Character traits and living a good life
				\item Legal compliance
			\end{enumerate}
		\end{block}
		\begin{exampleblock}{Answer: C}\end{exampleblock}
		\begin{block}{Explanation}
			Virtue ethics asks what kind of person one should be, emphasizing virtues like wisdom, courage, and justice.
		\end{block}
	\end{frame}
	
	\begin{frame}{S147 --- Contextual Integrity}
		\begin{block}{Question}
			Contextual integrity requires that:
			\begin{enumerate}
				\item Data can be freely reused
				\item Data use must respect the norms of the context in which it was provided
				\item Data must be encrypted
				\item Data must be deleted
			\end{enumerate}
		\end{block}
		\begin{exampleblock}{Answer: B}\end{exampleblock}
		\begin{block}{Explanation}
			Contextual integrity (Nissenbaum): data provided in one context should not be used in an incompatible context.
		\end{block}
	\end{frame}
	
	\begin{frame}{S148 --- Data Minimisation}
		\begin{block}{Question}
			Data minimisation requires:
			\begin{enumerate}
				\item Collecting all available data
				\item Collecting only data necessary for the specified purpose
				\item Sharing data freely
				\item Retaining data indefinitely
			\end{enumerate}
		\end{block}
		\begin{exampleblock}{Answer: B}\end{exampleblock}
		\begin{block}{Explanation}
			Data minimisation = collect only what is necessary for the specified purpose. Core GDPR principle.
		\end{block}
	\end{frame}
	
	\begin{frame}{S149 --- Right to be Forgotten}
		\begin{block}{Question}
			The right to be forgotten allows data subjects to:
			\begin{enumerate}
				\item Access their data
				\item Request erasure of personal data under certain conditions
				\item Sell their data
				\item Share data without consent
			\end{enumerate}
		\end{block}
		\begin{exampleblock}{Answer: B}\end{exampleblock}
		\begin{block}{Explanation}
			Right to erasure (GDPR Art. 17): allows deletion when data is no longer necessary, consent is withdrawn, etc.
		\end{block}
	\end{frame}
	
	\begin{frame}{S150 --- GDPR Breach Notification}
		\begin{block}{Question}
			GDPR requires notification of a personal data breach to the supervisory authority within:
			\begin{enumerate}
				\item 24 hours
				\item 48 hours
				\item 72 hours
				\item 1 week
			\end{enumerate}
		\end{block}
		\begin{exampleblock}{Answer: C}\end{exampleblock}
		\begin{block}{Explanation}
			GDPR Article 33 requires breach notification within 72 hours of becoming aware of the breach.
		\end{block}
	\end{frame}
	
	\begin{frame}{S151 --- GDPR Article 17}
		\begin{block}{Question}
			The GDPR right to erasure (right to be forgotten) is in:
			\begin{enumerate}
				\item Article 6
				\item Article 17
				\item Article 33
				\item Article 5
			\end{enumerate}
		\end{block}
		\begin{exampleblock}{Answer: B}\end{exampleblock}
		\begin{block}{Explanation}
			Article 17 = right to erasure. Article 6 = lawful basis; Article 33 = breach notification; Article 5 = principles.
		\end{block}
	\end{frame}
	
	\begin{frame}{S152 --- NIST AI RMF Functions}
		\begin{block}{Question}
			The NIST AI Risk Management Framework is organized around:
			\begin{enumerate}
				\item Train, Test, Deploy, Monitor
				\item Govern, Map, Measure, Manage
				\item Design, Build, Audit, Retire
				\item Plan, Do, Check, Act
			\end{enumerate}
		\end{block}
		\begin{exampleblock}{Answer: B}\end{exampleblock}
		\begin{block}{Explanation}
			NIST AI RMF = Govern (structures), Map (context), Measure (assessment), Manage (mitigation).
		\end{block}
	\end{frame}
	
	\begin{frame}{S153 --- EU AI Act Prohibited Systems}
		\begin{block}{Question}
			Which of the following is prohibited under the EU AI Act?
			\begin{enumerate}
				\item AI for medical diagnosis
				\item Social scoring based on social behavior
				\item AI for credit scoring
				\item AI for customer service
			\end{enumerate}
		\end{block}
		\begin{exampleblock}{Answer: B}\end{exampleblock}
		\begin{block}{Explanation}
			Prohibited: social scoring, biometric categorization using sensitive characteristics, emotion recognition in workplace/education, untargeted facial scraping.
		\end{block}
	\end{frame}
	
	\begin{frame}{S154 --- EU AI Act Fines}
		\begin{block}{Question}
			Non-compliance with the EU AI Act can lead to fines up to:
			\begin{enumerate}
				\item €10 million or 2\% of global turnover
				\item €35 million or 7\% of global turnover
				\item €5 million or 1\% of global turnover
				\item €100 million or 10\% of global turnover
			\end{enumerate}
		\end{block}
		\begin{exampleblock}{Answer: B}\end{exampleblock}
		\begin{block}{Explanation}
			EU AI Act fines range from €7.5M (1.5\%) to €35M (7\%) of global turnover depending on offense and company size.
		\end{block}
	\end{frame}
	
	\begin{frame}{S155 --- SR 11-7 Issuers}
		\begin{block}{Question}
			SR 11-7, ``Guidance on Model Risk Management'', was issued by:
			\begin{enumerate}
				\item The Federal Reserve and the OCC
				\item The SEC
				\item The FDIC alone
				\item The CFTC
			\end{enumerate}
		\end{block}
		\begin{exampleblock}{Answer: A}\end{exampleblock}
		\begin{block}{Explanation}
			SR 11-7 (2011) was issued jointly by the Federal Reserve and the Office of the Comptroller of the Currency.
		\end{block}
	\end{frame}
	
	\begin{frame}{S156 --- Three Lines of Defense}
		\begin{block}{Question}
			The Three Lines of Defense model was formalized by:
			\begin{enumerate}
				\item The Institute of Internal Auditors (IIA)
				\item The Federal Reserve alone
				\item The Basel Committee alone
				\item The SEC alone
			\end{enumerate}
		\end{block}
		\begin{exampleblock}{Answer: A}\end{exampleblock}
		\begin{block}{Explanation}
			The IIA formalized the Three Lines of Defense model, later updated to the Three Lines Model in 2020.
		\end{block}
	\end{frame}
	
	\begin{frame}{S157 --- Model Inventory Definition}
		\begin{block}{Question}
			A model inventory is:
			\begin{enumerate}
				\item A list of all models in use, with owners, risk tiers, and validation status
				\item A list of data sources
				\item A list of employees
				\item A list of IT systems
			\end{enumerate}
		\end{block}
		\begin{exampleblock}{Answer: A}\end{exampleblock}
		\begin{block}{Explanation}
			Model inventory = registry of all models. Includes owner, purpose, risk tier, validation status, and interdependencies.
		\end{block}
	\end{frame}
	
	\begin{frame}{S158 --- Model Landscape Definition}
		\begin{block}{Question}
			A model landscape is:
			\begin{enumerate}
				\item A list of all models
				\item A map of interactions and interdependencies between models
				\item A training dataset
				\item A model architecture
			\end{enumerate}
		\end{block}
		\begin{exampleblock}{Answer: B}\end{exampleblock}
		\begin{block}{Explanation}
			Model landscape = map of interactions and dependencies. It reveals cascade risk and critical nodes.
		\end{block}
	\end{frame}
	
	\begin{frame}{S159 --- Model Risk Appetite}
		\begin{block}{Question}
			Model risk appetite is:
			\begin{enumerate}
				\item The set of models in production
				\item The level of model risk the organization is willing to accept
				\item The cost of modeling
				\item The number of models
			\end{enumerate}
		\end{block}
		\begin{exampleblock}{Answer: B}\end{exampleblock}
		\begin{block}{Explanation}
			Model risk appetite = acceptable level of model risk, set by the board, guiding governance.
		\end{block}
	\end{frame}
	
	\begin{frame}{S160 --- Use Test Definition}
		\begin{block}{Question}
			The use test in model validation is:
			\begin{enumerate}
				\item A quantitative accuracy metric
				\item A qualitative assessment of whether the model is used appropriately for its intended purpose
				\item A code review
				\item A speed test
			\end{enumerate}
		\end{block}
		\begin{exampleblock}{Answer: B}\end{exampleblock}
		\begin{block}{Explanation}
			The use test assesses whether the model is used correctly and appropriately in practice. It is qualitative and people-focused.
		\end{block}
	\end{frame}
	
	\begin{frame}{S161 --- Outcomes Analysis}
		\begin{block}{Question}
			Outcomes analysis in model validation is used to:
			\begin{enumerate}
				\item Speed up training
				\item Compare model predictions against realized outcomes to detect bias, degradation, or incorrect assumptions
				\item Reduce features
				\item Increase accuracy
			\end{enumerate}
		\end{block}
		\begin{exampleblock}{Answer: B}\end{exampleblock}
		\begin{block}{Explanation}
			Outcomes analysis compares model predictions against actual outcomes to validate model quality in production.
		\end{block}
	\end{frame}
	
	\begin{frame}{S162 --- Initial Validation Goals}
		\begin{block}{Question}
			Initial validation serves which three key goals?
			\begin{enumerate}
				\item Operational feasibility, proper documentation, user training
				\item Accuracy, speed, cost
				\item Data collection, feature selection, model choice
				\item Deployment, monitoring, retirement
			\end{enumerate}
		\end{block}
		\begin{exampleblock}{Answer: A}\end{exampleblock}
		\begin{block}{Explanation}
			Initial validation ensures: (1) operational feasibility, (2) proper documentation and firm-wide standards, (3) user training.
		\end{block}
	\end{frame}
	
	\begin{frame}{S163 --- Effective Challenge}
		\begin{block}{Question}
			``Effective challenge'' in model risk management refers to:
			\begin{enumerate}
				\item A critical analysis by objective and informed parties to identify model assumptions and limitations
				\item A quantitative test
				\item A code review
				\item A data quality check
			\end{enumerate}
		\end{block}
		\begin{exampleblock}{Answer: A}\end{exampleblock}
		\begin{block}{Explanation}
			SR 11-7's ``effective challenge'': critical analysis by independent parties to identify assumptions, limitations, and weaknesses.
		\end{block}
	\end{frame}
	
	\begin{frame}{S164 --- Model Validation Frequency}
		\begin{block}{Question}
			According to SR 11-7, banks should conduct periodic review of each model at least:
			\begin{enumerate}
				\item Monthly
				\item Quarterly
				\item Annually, or more frequently if warranted
				\item Every 3 years
			\end{enumerate}
		\end{block}
		\begin{exampleblock}{Answer: C}\end{exampleblock}
		\begin{block}{Explanation}
			SR 11-7 recommends annual review at minimum, with more frequent reviews for higher-risk or AI/ML models.
		\end{block}
	\end{frame}
	
	\begin{frame}{S165 --- AI/ML Model Validation Differences}
		\begin{block}{Question}
			Compared to traditional models, AI/ML model validation requires:
			\begin{enumerate}
				\item Less frequent monitoring
				\item More frequent monitoring and broader stakeholder involvement
				\item No validation
				\item Only code review
			\end{enumerate}
		\end{block}
		\begin{exampleblock}{Answer: B}\end{exampleblock}
		\begin{block}{Explanation}
			AI/ML models require more frequent monitoring due to complexity, data drift, and broader stakeholder involvement (HR, compliance, operational risk).
		\end{block}
	\end{frame}
	
	\begin{frame}{S166 --- Model Risk Tier}
		\begin{block}{Question}
			The model risk tier determines:
			\begin{enumerate}
				\item The model's accuracy
				\item The depth and frequency of validation activities
				\item The number of features
				\item The training time
			\end{enumerate}
		\end{block}
		\begin{exampleblock}{Answer: B}\end{exampleblock}
		\begin{block}{Explanation}
			Risk tier drives validation depth and frequency. Higher-risk models are validated more frequently.
		\end{block}
	\end{frame}
	
	\begin{frame}{S167 --- Model Risk Ranking Factors}
		\begin{block}{Question}
			Model risk ranking typically includes which factors?
			\begin{enumerate}
				\item Materiality, complexity, exposure
				\item Accuracy, speed, cost
				\item Data size, feature count, training time
				\item Color, shape, size
			\end{enumerate}
		\end{block}
		\begin{exampleblock}{Answer: A}\end{exampleblock}
		\begin{block}{Explanation}
			Model risk ranking: materiality (impact), complexity (methodology), exposure (usage). These drive resource allocation.
		\end{block}
	\end{frame}
	
	\begin{frame}{S168 --- Data Provenance Definition}
		\begin{block}{Question}
			Data provenance refers to:
			\begin{enumerate}
				\item Data storage format
				\item The origin, history, and legal right to use data
				\item Data encryption
				\item Data classification
			\end{enumerate}
		\end{block}
		\begin{exampleblock}{Answer: B}\end{exampleblock}
		\begin{block}{Explanation}
			Data provenance = origin and legal rights to use data. Critical for alternative data and AI training sets.
		\end{block}
	\end{frame}
	
	\begin{frame}{S169 --- Data Classification Definition}
		\begin{block}{Question}
			Data classification is the process of:
			\begin{enumerate}
				\item Sorting data alphabetically
				\item Categorizing data by sensitivity and structure to determine handling and controls
				\item Encrypting data
				\item Compressing data
			\end{enumerate}
		\end{block}
		\begin{exampleblock}{Answer: B}\end{exampleblock}
		\begin{block}{Explanation}
			Data classification assigns sensitivity labels (public, internal, confidential, restricted) and structure labels.
		\end{block}
	\end{frame}
	
	\begin{frame}{S170 --- Metadata Definition}
		\begin{block}{Question}
			Metadata is best defined as:
			\begin{enumerate}
				\item Raw data
				\item Data about data
				\item Encrypted data
				\item Synthetic data
			\end{enumerate}
		\end{block}
		\begin{exampleblock}{Answer: B}\end{exampleblock}
		\begin{block}{Explanation}
			Metadata = data about data. It documents origin, structure, definitions, and relationships.
		\end{block}
	\end{frame}
	
	\begin{frame}{S171 --- RBAC Definition}
		\begin{block}{Question}
			Role-Based Access Control (RBAC) grants access based on:
			\begin{enumerate}
				\item Individual discretion
				\item User roles within the organization
				\item Random selection
				\item Data size
			\end{enumerate}
		\end{block}
		\begin{exampleblock}{Answer: B}\end{exampleblock}
		\begin{block}{Explanation}
			RBAC = permissions based on organizational roles. ABAC = attribute-based; MAC = mandatory; DAC = discretionary.
		\end{block}
	\end{frame}
	
	\begin{frame}{S172 --- ABAC Definition}
		\begin{block}{Question}
			Attribute-Based Access Control (ABAC) grants access based on:
			\begin{enumerate}
				\item User roles
				\item Attributes of user, resource, and environment
				\item Random selection
				\item Seniority
			\end{enumerate}
		\end{block}
		\begin{exampleblock}{Answer: B}\end{exampleblock}
		\begin{block}{Explanation}
			ABAC uses attributes (user, resource, environment) to make access decisions. More flexible than RBAC.
		\end{block}
	\end{frame}
	
	\begin{frame}{S173 --- Data Strategy Elements}
		\begin{block}{Question}
			A data strategy should include:
			\begin{enumerate}
				\item Vision, goals, priorities, and data governance policy
				\item Only data storage
				\item Only data encryption
				\item Only data deletion
			\end{enumerate}
		\end{block}
		\begin{exampleblock}{Answer: A}\end{exampleblock}
		\begin{block}{Explanation}
			Data strategy includes vision, goals, priorities, authorization for non-public data, and approved data-governance policy.
		\end{block}
	\end{frame}
	
	\begin{frame}{S174 --- Data Governance Committee}
		\begin{block}{Question}
			A data governance committee is typically responsible for:
			\begin{enumerate}
				\item Operational data entry
				\item Approved data practices and policies, including collection, treatment, access, and monitoring
				\item Only data deletion
				\item Only data encryption
			\end{enumerate}
		\end{block}
		\begin{exampleblock}{Answer: B}\end{exampleblock}
		\begin{block}{Explanation}
			Data governance committee provides oversight on all aspects of data governance (collection, reconciliations, treatment, access, monitoring).
		\end{block}
	\end{frame}
	
	\begin{frame}{S175 --- Synthetic Data Purpose}
		\begin{block}{Question}
			Synthetic data is used to:
			\begin{enumerate}
				\item Replace real data entirely
				\item Augment available data and mitigate privacy risks
				\item Increase data storage
				\item Slow down training
			\end{enumerate}
		\end{block}
		\begin{exampleblock}{Answer: B}\end{exampleblock}
		\begin{block}{Explanation}
			Synthetic data augments training data and mitigates risks when personally identifiable data may be present.
		\end{block}
	\end{frame}
	
	\begin{frame}{S176 --- AI-Ready Data}
		\begin{block}{Question}
			AI-ready data is:
			\begin{enumerate}
				\item Raw and unstructured
				\item Accurate, unbiased, well-structured, cleaned, and suitable for training
				\item Encrypted only
				\item Deleted data
			\end{enumerate}
		\end{block}
		\begin{exampleblock}{Answer: B}\end{exampleblock}
		\begin{block}{Explanation}
			AI-ready data is accurate, unbiased, well-structured, cleaned, and suitable for training AI models.
		\end{block}
	\end{frame}
	
	\begin{frame}{S177 --- Data Retention Policy}
		\begin{block}{Question}
			A data retention policy specifies:
			\begin{enumerate}
				\item How long data is kept and when it is deleted
				\item Only data storage format
				\item Only data encryption
				\item Only data access
			\end{enumerate}
		\end{block}
		\begin{exampleblock}{Answer: A}\end{exampleblock}
		\begin{block}{Explanation}
			Data retention policy defines how long data is kept and procedures for deletion, complying with regulations.
		\end{block}
	\end{frame}
	
	\begin{frame}{S178 --- GDPR Data Subject Rights}
		\begin{block}{Question}
			Under GDPR, data subjects have the right to:
			\begin{enumerate}
				\item Only access data
				\item Access, rectify, erase, restrict processing, data portability, and object
				\item Only delete data
				\item Only sell data
			\end{enumerate}
		\end{block}
		\begin{exampleblock}{Answer: B}\end{exampleblock}
		\begin{block}{Explanation}
			GDPR rights: access, rectification, erasure, restriction, portability, objection, and rights related to automated decision-making.
		\end{block}
	\end{frame}
	
	\begin{frame}{S179 --- GDPR Extraterritorial Scope}
		\begin{block}{Question}
			GDPR applies to organizations outside the EEA if they:
			\begin{enumerate}
				\item Have any European customers
				\item Offer services to data subjects within the EEA or monitor their behavior
				\item Are publicly traded
				\item Have more than 100 employees
			\end{enumerate}
		\end{block}
		\begin{exampleblock}{Answer: B}\end{exampleblock}
		\begin{block}{Explanation}
			GDPR has extraterritorial jurisdiction: applies to organizations outside EEA offering services to or monitoring EEA data subjects.
		\end{block}
	\end{frame}
	
	\begin{frame}{S180 --- CCPA Full Name}
		\begin{block}{Question}
			CCPA stands for:
			\begin{enumerate}
				\item California Consumer Privacy Act
				\item California Cyber Protection Act
				\item California Consumer Protection Act
				\item California Data Privacy Act
			\end{enumerate}
		\end{block}
		\begin{exampleblock}{Answer: A}\end{exampleblock}
		\begin{block}{Explanation}
			CCPA = California Consumer Privacy Act (2018). Most comprehensive US state privacy law.
		\end{enumerate}
	\end{block}
\end{frame}

\begin{frame}{S181 --- CCPA Consumer Rights}
	\begin{block}{Question}
		Under CCPA, consumers have the right to:
		\begin{enumerate}
			\item Know, access, delete, opt out of sale, and non-discrimination
			\item Only access data
			\item Only delete data
			\item Only sell data
		\end{enumerate}
	\end{block}
	\begin{exampleblock}{Answer: A}\end{exampleblock}
	\begin{block}{Explanation}
		CCPA rights: know what data is collected, access, delete, opt out of sale, and non-discrimination.
	\end{block}
\end{frame}

\begin{frame}{S182 --- BIPA Full Name}
	\begin{block}{Question}
		BIPA stands for:
		\begin{enumerate}
			\item Biometric Information Privacy Act
			\item Business Information Protection Act
			\item Biometric Identification Protection Act
			\item Background Information Privacy Act
		\end{enumerate}
	\end{block}
	\begin{exampleblock}{Answer: A}\end{exampleblock}
	\begin{block}{Explanation}
		BIPA = Biometric Information Privacy Act (Illinois, 2008). Regulates collection of biometric identifiers.
	\end{block}
\end{frame}

\begin{frame}{S183 --- BIPA Statutory Damages}
	\begin{block}{Question}
		Under BIPA, statutory damages are:
		\begin{enumerate}
			\item \$100 per violation
			\item \$1,000 per violation (\$5,000 if intentional/reckless)
			\item \$10,000 per violation
			\item No damages
		\end{enumerate}
	\end{block}
	\begin{exampleblock}{Answer: B}\end{exampleblock}
	\begin{block}{Explanation}
		BIPA: \$1,000 per violation, or \$5,000 if intentional or reckless. Individuals can sue for damages.
	\end{block}
\end{frame}

\begin{frame}{S184 --- HIPAA Full Name}
	\begin{block}{Question}
		HIPAA stands for:
		\begin{enumerate}
			\item Health Insurance Portability and Accountability Act
			\item Health Information Privacy and Accountability Act
			\item Health Insurance Privacy and Protection Act
			\item Health Information Portability Act
		\end{enumerate}
	\end{block}
	\begin{exampleblock}{Answer: A}\end{exampleblock}
	\begin{block}{Explanation}
		HIPAA = Health Insurance Portability and Accountability Act (1996). Protects privacy and security of health information.
	\end{block}
\end{frame}

\begin{frame}{S185 --- GLBA Full Name}
	\begin{block}{Question}
		GLBA stands for:
		\begin{enumerate}
			\item Gramm-Leach-Bliley Act
			\item General Liability and Business Act
			\item Government-Led Banking Act
			\item Global Banking and Lending Act
		\end{enumerate}
	\end{block}
	\begin{exampleblock}{Answer: A}\end{exampleblock}
	\begin{block}{Explanation}
		GLBA = Gramm-Leach-Bliley Act. Requires financial institutions to explain privacy practices and implement security safeguards.
	\end{block}
\end{frame}

\begin{frame}{S186 --- FTC Safeguards Rule}
	\begin{block}{Question}
		The FTC Safeguards Rule requires:
		\begin{enumerate}
			\item Financial institutions to have measures to keep customer information secure
			\item Only data deletion
			\item Only data encryption
			\item Only data access
		\end{enumerate}
	\end{block}
	\begin{exampleblock}{Answer: A}\end{exampleblock}
	\begin{block}{Explanation}
		FTC Safeguards Rule: financial institutions under FTC jurisdiction must have measures to keep customer information secure.
	\end{block}
\end{frame}

\begin{frame}{S187 --- Privacy Act of 1974}
	\begin{block}{Question}
		The Privacy Act of 1974 governs:
		\begin{enumerate}
			\item How federal agencies collect and use data about individuals
			\item Private sector data collection
			\item Only health data
			\item Only financial data
		\end{enumerate}
	\end{block}
	\begin{exampleblock}{Answer: A}\end{exampleblock}
	\begin{block}{Explanation}
		Privacy Act of 1974: governs federal agencies' collection and use of personal data; prohibits disclosure without consent.
	\end{block}
\end{frame}

\begin{frame}{S188 --- EU Digital Markets Act}
	\begin{block}{Question}
		The EU Digital Markets Act (DMA) targets:
		\begin{enumerate}
			\item All companies
			\item Large online platforms with ``gatekeeper'' status
			\item Only banks
			\item Only healthcare companies
		\end{enumerate}
	\end{block}
	\begin{exampleblock}{Answer: B}\end{exampleblock}
	\begin{block}{Explanation}
		DMA targets ``gatekeepers'' -- large platforms with dominant positions. Bans self-preferencing, requires interoperability.
	\end{block}
\end{frame}

\begin{frame}{S189 --- EU Digital Services Act}
	\begin{block}{Question}
		The EU Digital Services Act (DSA) requires platforms to:
		\begin{enumerate}
			\item Only remove illegal content
			\item Prevent illegal content, be transparent about moderation, address disinformation, protect against algorithmic bias
			\item Only label deepfakes
			\item Only allow opt-out
		\end{enumerate}
	\end{block}
	\begin{exampleblock}{Answer: B}\end{exampleblock}
	\begin{block}{Explanation}
		DSA: prevent illegal content, transparent moderation, address disinformation, protect against algorithmic bias, allow opt-out.
	\end{block}
\end{frame}

\begin{frame}{S190 --- China Generative AI Measures}
	\begin{block}{Question}
		China's Generative AI Measures require service providers to:
		\begin{enumerate}
			\item Only register with government
			\item Prohibit illegal content, prevent discriminatory content, not infringe on privacy rights
			\item Only label AI content
			\item Only obtain consent
		\end{enumerate}
	\end{block}
	\begin{exampleblock}{Answer: B}\end{exampleblock}
	\begin{block}{Explanation}
		China's Generative AI Measures (2023): prohibit illegal content, prevent discrimination, protect privacy and personal information.
	\end{block}
\end{frame}

\begin{frame}{S191 --- PIPL Full Name}
	\begin{block}{Question}
		PIPL stands for:
		\begin{enumerate}
			\item Personal Information Protection Law
			\item Privacy and Information Protection Law
			\item Personal Identification Protection Law
			\item Public Information Privacy Law
		\end{enumerate}
	\end{block}
	\begin{exampleblock}{Answer: A}\end{exampleblock}
	\begin{block}{Explanation}
		PIPL = Personal Information Protection Law (China). Mandates data minimization, user consent, and algorithmic transparency.
	\end{block}
\end{frame}

\begin{frame}{S192 --- Colorado Privacy Act}
	\begin{block}{Question}
		The Colorado Privacy Act (CPA) gives consumers the right to:
		\begin{enumerate}
			\item Only access data
			\item Know, access, delete, opt out of sale, and opt out of targeted advertising
			\item Only delete data
			\item Only sell data
		\end{enumerate}
	\end{block}
	\begin{exampleblock}{Answer: B}\end{exampleblock}
	\begin{block}{Explanation}
		CPA (2021): know, access, delete, opt out of sale, and opt out of targeted advertising. Includes data minimization provisions.
	\end{block}
\end{frame}

\begin{frame}{S193 --- Virginia CDPA}
	\begin{block}{Question}
		The Virginia Consumer Data Protection Act (CDPA) applies to businesses that:
		\begin{enumerate}
			\item Have any Virginia customers
			\item Control or process personal data of at least 100,000 Virginia consumers
			\item Are publicly traded
			\item Have more than 10 employees
		\end{enumerate}
	\end{block}
	\begin{exampleblock}{Answer: B}\end{exampleblock}
	\begin{block}{Explanation}
		CDPA applies to businesses controlling/processing data of ≥100,000 Virginia consumers, or ≥25,000 with >50\% revenue from data sales.
	\end{block}
\end{frame}

\begin{frame}{S194 --- Illinois PIPA}
	\begin{block}{Question}
		The Illinois Personal Information Protection Act (PIPA) requires:
		\begin{enumerate}
			\item Only data encryption
			\item Reasonable security measures, breach notification within 24 hours, privacy notice, access, and opt-out
			\item Only data deletion
			\item Only data access
		\end{enumerate}
	\end{block}
	\begin{exampleblock}{Answer: B}\end{exampleblock}
	\begin{block}{Explanation}
		PIPA: reasonable security, breach notification within 24 hours, privacy notice, access, correction, and opt-out of sale.
	\end{block}
\end{frame}

\begin{frame}{S195 --- EU AI Act Effective Date}
	\begin{block}{Question}
		The EU AI Act became effective on:
		\begin{enumerate}
			\item August 1, 2024
			\item January 1, 2025
			\item March 13, 2024
			\item May 21, 2024
		\end{enumerate}
	\end{block}
	\begin{exampleblock}{Answer: A}\end{exampleblock}
	\begin{block}{Explanation}
		EU AI Act effective August 1, 2024. Ban on unacceptable risk systems from February 2, 2025. Most provisions from August 2, 2026.
	\end{block}
\end{frame}

\begin{frame}{S196 --- European AI Office}
	\begin{block}{Question}
		The European AI Office is responsible for:
		\begin{enumerate}
			\item Only research
			\item Compliance, implementation, and enforcement of the EU AI Act
			\item Only data protection
			\item Only cybersecurity
		\end{enumerate}
	\end{block}
	\begin{exampleblock}{Answer: B}\end{exampleblock}
	\begin{block}{Explanation}
		European AI Office: first body in the world to enforce binding rules on AI. Responsible for compliance, implementation, enforcement.
	\end{block}
\end{frame}

\begin{frame}{S197 --- AI Act High-Risk Requirements}
	\begin{block}{Question}
		Providers of high-risk AI under the EU AI Act must:
		\begin{enumerate}
			\item Only register
			\item Keep thorough records of datasets, methodologies, and oversight measures
			\item Only label content
			\item Only obtain consent
		\end{enumerate}
	\end{block}
	\begin{exampleblock}{Answer: B}\end{exampleblock}
	\begin{block}{Explanation}
		High-risk AI providers must keep records of datasets, programming/training methodologies, and oversight measures. Impact assessment required.
	\end{block}
\end{frame}

\begin{frame}{S198 --- AI Act Prohibited: Biometric Categorization}
	\begin{block}{Question}
		Under the EU AI Act, biometric categorization systems using sensitive characteristics are:
		\begin{enumerate}
			\item Allowed with consent
			\item Prohibited
			\item Allowed for security
			\item Allowed for research
		\end{enumerate}
	\end{block}
	\begin{exampleblock}{Answer: B}\end{exampleblock}
	\begin{block}{Explanation}
		Prohibited: biometric categorization using political, religious, philosophical beliefs, sexual orientation, race.
	\end{block}
\end{frame}

\begin{frame}{S199 --- AI Act Prohibited: Emotion Recognition}
	\begin{block}{Question}
		Under the EU AI Act, emotion recognition in the workplace and educational institutions is:
		\begin{enumerate}
			\item Allowed
			\item Prohibited
			\item Allowed with consent
			\item Allowed for safety
		\end{enumerate}
	\end{block}
	\begin{exampleblock}{Answer: B}\end{exampleblock}
	\begin{block}{Explanation}
		Prohibited: emotion recognition in workplace and educational institutions under the EU AI Act.
	\end{block}
\end{frame}

\begin{frame}{S200 --- AI Act Prohibited: Social Scoring}
	\begin{block}{Question}
		Under the EU AI Act, social scoring based on social behavior or personal characteristics is:
		\begin{enumerate}
			\item Allowed
			\item Prohibited
			\item Allowed with consent
			\item Allowed for government
		\end{enumerate}
	\end{block}
	\begin{exampleblock}{Answer: B}\end{exampleblock}
	\begin{block}{Explanation}
		Prohibited: social scoring based on social behavior or personal characteristics under the EU AI Act.
	\end{block}
\end{frame}

\begin{frame}{S201 --- AI Act Prohibited: Manipulation}
	\begin{block}{Question}
		Under the EU AI Act, AI systems that manipulate human behavior to circumvent free will are:
		\begin{enumerate}
			\item Allowed
			\item Prohibited
			\item Allowed with consent
			\item Allowed for marketing
		\end{enumerate}
	\end{block}
	\begin{exampleblock}{Answer: B}\end{exampleblock}
	\begin{block}{Explanation}
		Prohibited: AI systems that manipulate human behavior to circumvent free will, and AI exploiting vulnerabilities.
	\end{block}
\end{frame}

\begin{frame}{S202 --- AI Act Prohibited: Facial Scraping}
	\begin{block}{Question}
		Under the EU AI Act, untargeted scraping of facial images to create facial recognition databases is:
		\begin{enumerate}
			\item Allowed
			\item Prohibited
			\item Allowed with consent
			\item Allowed for security
		\end{enumerate}
	\end{block}
	\begin{exampleblock}{Answer: B}\end{exampleblock}
	\begin{block}{Explanation}
		Prohibited: untargeted scraping of facial images from internet or CCTV to create facial recognition databases.
	\end{block}
\end{frame}

\begin{frame}{S203 --- AI Act Transparency Obligations}
	\begin{block}{Question}
		The EU AI Act requires companies to:
		\begin{enumerate}
			\item Only register
			\item Notify people when interacting with chatbots, biometric systems, or emotion recognition; label deepfakes
			\item Only label content
			\item Only obtain consent
		\end{enumerate}
	\end{block}
	\begin{exampleblock}{Answer: B}\end{exampleblock}
	\begin{block}{Explanation}
		AI Act transparency: notify when interacting with chatbots/biometric/emotion recognition; label deepfakes and AI-generated content.
	\end{block}
\end{frame}

\begin{frame}{S204 --- AI Act Impact Assessment}
	\begin{block}{Question}
		Under the EU AI Act, organizations deploying high-risk AI must:
		\begin{enumerate}
			\item Only register
			\item Do an impact assessment on how AI affects people's rights
			\item Only label content
			\item Only obtain consent
		\end{enumerate}
	\end{block}
	\begin{exampleblock}{Answer: B}\end{exampleblock}
	\begin{block}{Explanation}
		High-risk AI deployers (banks, insurance, essential services) must do an impact assessment on fundamental rights.
	\end{block}
\end{frame}

\begin{frame}{S205 --- AI Act Fines Range}
	\begin{block}{Question}
		EU AI Act fines range from:
		\begin{enumerate}
			\item €1M to €10M
			\item €7.5M (1.5\%) to €35M (7\%) of global turnover
			\item €100M to €500M
			\item No fines
		\end{enumerate}
	\end{block}
	\begin{exampleblock}{Answer: B}\end{exampleblock}
	\begin{block}{Explanation}
		EU AI Act fines: €7.5M or 1.5\% to €35M or 7\% of global turnover, depending on offense and company size.
	\end{block}
\end{frame}

\begin{frame}{S206 --- Executive Order on AI (US)}
	\begin{block}{Question}
		President Biden's Executive Order on Safe, Secure, and Trustworthy AI was issued on:
		\begin{enumerate}
			\item October 30, 2023
			\item January 1, 2024
			\item May 12, 2021
			\item August 1, 2024
		\end{enumerate}
	\end{block}
	\begin{exampleblock}{Answer: A}\end{exampleblock}
	\begin{block}{Explanation}
		EO on Safe, Secure, and Trustworthy Development and Use of AI issued October 30, 2023. Requires safety test results sharing with government.
	\end{block}
\end{frame}

\begin{frame}{S207 --- AI Safety and Security Board}
	\begin{block}{Question}
		The AI Safety and Security Board was established by:
		\begin{enumerate}
			\item The White House
			\item The Department of Homeland Security
			\item The SEC
			\item The FTC
		\end{enumerate}
	\end{block}
	\begin{exampleblock}{Answer: B}\end{exampleblock}
	\begin{block}{Explanation}
		AI Safety and Security Board established by DHS on April 26, 2024, to advise on safe AI development in critical infrastructure.
	\end{block}
\end{frame}

\begin{frame}{S208 --- NIST AI RMF Functions}
	\begin{block}{Question}
		The NIST AI RMF's four functions are:
		\begin{enumerate}
			\item Govern, Map, Measure, Manage
			\item Plan, Do, Check, Act
			\item Train, Test, Deploy, Monitor
			\item Design, Build, Audit, Retire
		\end{enumerate}
	\end{block}
	\begin{exampleblock}{Answer: A}\end{exampleblock}
	\begin{block}{Explanation}
		NIST AI RMF: Govern (structures), Map (context), Measure (assessment), Manage (mitigation). Voluntary framework.
	\end{block}
\end{frame}

\begin{frame}{S209 --- NIST AI RMF Nature}
	\begin{block}{Question}
		The NIST AI Risk Management Framework is:
		\begin{enumerate}
			\item Mandatory
			\item Voluntary and non-prescriptive
			\item Only for government
			\item Only for banks
		\end{enumerate}
	\end{block}
	\begin{exampleblock}{Answer: B}\end{exampleblock}
	\begin{block}{Explanation}
		NIST AI RMF is voluntary, non-prescriptive, flexible, and customizable based on specific use cases and risk profiles.
	\end{block}
\end{frame}

\begin{frame}{S210 --- SR 11-7 Definition of Model}
	\begin{block}{Question}
		SR 11-7 defines a model as:
		\begin{enumerate}
			\item Any spreadsheet
			\item A quantitative method, system, or approach that applies statistical, economic, financial, or mathematical theories to process input data into quantitative estimates
			\item Only AI models
			\item Only machine learning models
		\end{enumerate}
	\end{block}
	\begin{exampleblock}{Answer: B}\end{exampleblock}
	\begin{block}{Explanation}
		SR 11-7: model = quantitative method/system/approach applying theories and assumptions to process input data into quantitative estimates.
	\end{block}
\end{frame}

\begin{frame}{S211 --- Model Risk Definition}
	\begin{block}{Question}
		Model risk is defined as:
		\begin{enumerate}
			\item The risk of data loss
			\item The risk of adverse consequences from decisions based on incorrect or misused model outputs
			\item The risk of cyber attack
			\item The risk of model theft
		\end{enumerate}
	\end{block}
	\begin{exampleblock}{Answer: B}\end{exampleblock}
	\begin{block}{Explanation}
		Model risk = risk of adverse consequences from decisions based on incorrect or misused model outputs and reports.
	\end{block}
\end{frame}

\begin{frame}{S212 --- Model Inventory Contents}
	\begin{block}{Question}
		A comprehensive model inventory should include:
		\begin{enumerate}
			\item Only model names
			\item Owner, purpose, methodology, risk tier, validation status, interdependencies
			\item Only model code
			\item Only model data
		\end{enumerate}
	\end{block}
	\begin{exampleblock}{Answer: B}\end{exampleblock}
	\begin{block}{Explanation}
		Model inventory includes: owner, purpose, methodology, risk classification, validation status, interdependencies, and more.
	\end{block}
\end{frame}

\begin{frame}{S213 --- Model Owner Responsibility}
	\begin{block}{Question}
		The model owner is responsible for:
		\begin{enumerate}
			\item Only model development
			\item Proper development, validation, approval, remediation, implementation, use, monitoring, and change management
			\item Only model use
			\item Only model documentation
		\end{enumerate}
	\end{block}
	\begin{exampleblock}{Answer: B}\end{exampleblock}
	\begin{block}{Explanation}
		Model owner oversees the model lifecycle: development, validation, approval, findings remediation, implementation, use, monitoring, and change management.
	\end{block}
\end{frame}

\begin{frame}{S214 --- Model Developer Responsibility}
	\begin{block}{Question}
		Model developers are responsible for:
		\begin{enumerate}
			\item Only coding
			\item Creating models, documenting decisions, and responding to validation feedback
			\item Only testing
			\item Only deployment
		\end{enumerate}
	\end{block}
	\begin{exampleblock}{Answer: B}\end{exampleblock}
	\begin{block}{Explanation}
		Model developers create models, document modeling decisions, and act on validation feedback to improve models.
	\end{block}
\end{frame}

\begin{frame}{S215 --- Model Validator Responsibility}
	\begin{block}{Question}
		Model validators are responsible for:
		\begin{enumerate}
			\item Developing models
			\item Providing independent and objective evaluation of models, assumptions, and methodologies
			\item Only coding
			\item Only data collection
		\end{enumerate}
	\end{block}
	\begin{exampleblock}{Answer: B}\end{exampleblock}
	\begin{block}{Explanation}
		Model validators provide independent and objective evaluation, reviewing assumptions, methodologies, and communicating findings to owners/developers.
	\end{block}
\end{frame}

\begin{frame}{S216 --- Internal Audit Role in Model Risk}
	\begin{block}{Question}
		Internal and external auditors in model risk management:
		\begin{enumerate}
			\item Are the primary model validators
			\item Provide independent checks on the validation process and outputs
			\item Develop models
			\item Own models
		\end{enumerate}
	\end{block}
	\begin{exampleblock}{Answer: B}\end{exampleblock}
	\begin{block}{Explanation}
		Auditors provide ``meta-validation'' -- independent checks on the validation process and outputs, not primary validation.
	\end{block}
\end{frame}

\begin{frame}{S217 --- Regulator Role in Model Risk}
	\begin{block}{Question}
		Regulatory authorities may review:
		\begin{enumerate}
			\item Only financial statements
			\item Risk models' validation processes and outcomes as part of supervision
			\item Only IT systems
			\item Only data
		\end{enumerate}
	\end{block}
	\begin{exampleblock}{Answer: B}\end{exampleblock}
	\begin{block}{Explanation}
		Regulators review risk model validation processes and outcomes, including internal audits and compliance with regulatory requirements.
	\end{block}
\end{frame}

\begin{frame}{S218 --- Model Risk Committee}
	\begin{block}{Question}
		A model risk committee is used to:
		\begin{enumerate}
			\item Only develop models
			\item Ensure comparable information levels among participants and shared responsibility for difficult decisions
			\item Only approve models
			\item Only validate models
		\end{enumerate}
	\end{block}
	\begin{exampleblock}{Answer: B}\end{exampleblock}
	\begin{block}{Explanation}
		Model risk committee ensures comparable information and shared responsibility for difficult decisions.
	\end{block}
\end{frame}

\begin{frame}{S219 --- Three Lines of Defense: First Line}
	\begin{block}{Question}
		The first line of defense in model risk management is:
		\begin{enumerate}
			\item Independent risk management
			\item Business and model owners/developers
			\item Internal audit
			\item Regulators
		\end{enumerate}
	\end{block}
	\begin{exampleblock}{Answer: B}\end{exampleblock}
	\begin{block}{Explanation}
		First line: business and model owners/developers. They own the model and its performance, monitor metrics, and react to breaches.
	\end{block}
\end{frame}

\begin{frame}{S220 --- Three Lines of Defense: Second Line}
	\begin{block}{Question}
		The second line of defense in model risk management is:
		\begin{enumerate}
			\item Business units
			\item Independent risk management/model validation
			\item Internal audit
			\item Regulators
		\end{enumerate}
	\end{block}
	\begin{exampleblock}{Answer: B}\end{exampleblock}
	\begin{block}{Explanation}
		Second line: independent risk management/model validation. Sets policy, conducts validations, challenges first line, maintains inventory.
	\end{block}
\end{frame}

\begin{frame}{S221 --- Three Lines of Defense: Third Line}
	\begin{block}{Question}
		The third line of defense in model risk management is:
		\begin{enumerate}
			\item Business units
			\item Independent risk management
			\item Internal audit
			\item Regulators
		\end{enumerate}
	\end{block}
	\begin{exampleblock}{Answer: C}\end{exampleblock}
	\begin{block}{Explanation}
		Third line: internal audit. Provides independent assurance to the board that the model risk management framework is well-designed and operating effectively.
	\end{block}
\end{frame}

\begin{frame}{S222 --- AI/ML Validation: Explainability Focus}
	\begin{block}{Question}
		For black-box AI/ML models, validation should focus on:
		\begin{enumerate}
			\item Only code review
			\item Explainability, benchmarking, assessment on different datasets, and back testing
			\item Only data quality
			\item Only model accuracy
		\end{enumerate}
	\end{block}
	\begin{exampleblock}{Answer: B}\end{exampleblock}
	\begin{block}{Explanation}
		Black-box validation: focus on explainability, benchmarking against classical models, assessment on different datasets, and back testing.
	\end{block}
\end{frame}

\begin{frame}{S223 --- AI/ML Validation: Benchmarking}
	\begin{block}{Question}
		Benchmarking in AI/ML model validation involves:
		\begin{enumerate}
			\item Only comparing accuracy
			\item Comparing model results to classical models to understand deviations
			\item Only data quality checks
			\item Only code review
		\end{enumerate}
	\end{block}
	\begin{exampleblock}{Answer: B}\end{exampleblock}
	\begin{block}{Explanation}
		Benchmarking: comparison of AI/ML model results to classical models to understand reasons for deviation.
	\end{block}
\end{frame}

\begin{frame}{S224 --- AI/ML Validation: In-Sample vs Holdout}
	\begin{block}{Question}
		SR 11-7's advice on analyzing ``in-sample fit and model performance in holdout samples'' is:
		\begin{enumerate}
			\item Not applicable to AI/ML
			\item One of the main applicable tests for AI/ML models
			\item Only for traditional models
			\item Only for vendor models
		\end{enumerate}
	\end{block}
	\begin{exampleblock}{Answer: B}\end{exampleblock}
	\begin{block}{Explanation}
		SR 11-7's holdout sample analysis is one of the main applicable tests for AI/ML models.
	\end{block}
\end{frame}

\begin{frame}{S225 --- AI/ML Validation: Back Testing}
	\begin{block}{Question}
		For AI/ML models, common back testing methods like MSE may become ineffective because:
		\begin{enumerate}
			\item MSE is always effective
			\item There may no longer be a straightforward relation between inputs and outputs
			\item MSE is too complex
			\item MSE is not available
		\end{enumerate}
	\end{block}
	\begin{exampleblock}{Answer: B}\end{exampleblock}
	\begin{block}{Explanation}
		For AI/ML, MSE may be ineffective due to non-straightforward input-output relations. K-fold cross-validation may be more suitable.
	\end{block}
\end{frame}

\begin{frame}{S226 --- AI/ML Validation: K-Fold CV}
	\begin{block}{Question}
		For AI/ML models, which technique might be more suitable for back testing than MSE?
		\begin{enumerate}
			\item K-fold cross-validation
			\item Only R-squared
			\item Only p-values
			\item Only t-statistics
		\end{enumerate}
	\end{block}
	\begin{exampleblock}{Answer: A}\end{exampleblock}
	\begin{block}{Explanation}
		K-fold cross-validation is more suitable for AI/ML back testing than traditional MSE due to non-linear input-output relations.
	\end{block}
\end{frame}

\begin{frame}{S227 --- AI/ML Validation: Reporting Component}
	\begin{block}{Question}
		The reporting component in AI/ML validation verifies that model outputs are:
		\begin{enumerate}
			\item Only accurate
			\item Accurate, complete, informative, and contain appropriate indicators of model performance and limitations
			\item Only fast
			\item Only cheap
		\end{enumerate}
	\end{block}
	\begin{exampleblock}{Answer: B}\end{exampleblock}
	\begin{block}{Explanation}
		Reporting component: verify outputs are accurate, complete, informative, and contain indicators of performance and limitations.
	\end{block}
\end{frame}

\begin{frame}{S228 --- AI/ML Validation: Specific Cases}
	\begin{block}{Question}
		Assessment of specific cases in AI/ML validation includes:
		\begin{enumerate}
			\item Only extreme input values
			\item Stress testing with extreme values and analysis of borderline decisions
			\item Only average cases
			\item Only random cases
		\end{enumerate}
	\end{block}
	\begin{exampleblock}{Answer: B}\end{exampleblock}
	\begin{block}{Explanation}
		Specific case assessment: stress testing with extreme input values and analysis of decisions made for borderline cases.
	\end{block}
\end{frame}

\begin{frame}{S229 --- AI/ML Validation: Focus on Explainability}
	\begin{block}{Question}
		For AI/ML models, explainability is key because:
		\begin{enumerate}
			\item Models are always transparent
			\item Models are often less than transparent and may be black-box
			\item Models are simple
			\item Models are always correct
		\end{enumerate}
	\end{block}
	\begin{exampleblock}{Answer: B}\end{exampleblock}
	\begin{block}{Explanation}
		AI/ML models may be less transparent (black-box). Explainability is key; outcomes must be rigorously tested and compared with traditional models.
	\end{block}
\end{frame}

\begin{frame}{S230 --- GenAI Validation Challenges}
	\begin{block}{Question}
		Which of the following is a challenge in validating GenAI models?
		\begin{enumerate}
			\item Deterministic output
			\item Open-ended output and hallucination
			\item Simple performance metrics
			\item No data drift
		\end{enumerate}
	\end{block}
	\begin{exampleblock}{Answer: B}\end{exampleblock}
	\begin{block}{Explanation}
		GenAI validation challenges: open-ended output, complex performance metrics, data drift, hallucination, bias, multi-modal issues, and interpretability.
	\end{block}
\end{frame}

\begin{frame}{S231 --- GenAI Hallucination Risk}
	\begin{block}{Question}
		Hallucination in GenAI models refers to:
		\begin{enumerate}
			\item Generating factually incorrect but convincing output
			\item Generating random noise
			\item Generating too many tokens
			\item Generating too few tokens
		\end{enumerate}
	\end{block}
	\begin{exampleblock}{Answer: A}\end{exampleblock}
	\begin{block}{Explanation}
		Hallucination: factually incorrect, misleading, or fabricated output that appears credible and convincing.
	\end{block}
\end{frame}

\begin{frame}{S232 --- GenAI Data Drift}
	\begin{block}{Question}
		Data drift in GenAI models refers to:
		\begin{enumerate}
			\item Changes in model architecture
			\item Changes in underlying environment making model responses obsolete
			\item Changes in training data only
			\item Changes in learning rate
		\end{enumerate}
	\end{block}
	\begin{exampleblock}{Answer: B}\end{exampleblock}
	\begin{block}{Explanation}
		Data drift: with passage of time, model responses may become irrelevant due to changes in underlying environment, making model obsolete.
	\end{block}
\end{frame}

\begin{frame}{S233 --- GenAI Bias Amplification}
	\begin{block}{Question}
		AI models can amplify bias inherent in data because:
		\begin{enumerate}
			\item They are unbiased
			\item They learn and reproduce patterns, including discriminatory ones
			\item They ignore data
			\item They only use clean data
		\end{enumerate}
	\end{block}
	\begin{exampleblock}{Answer: B}\end{exampleblock}
	\begin{block}{Explanation}
		AI models can amplify bias in data. Proper attention must ensure model responses are not discriminatory or biased.
	\end{block}
\end{frame}

\begin{frame}{S234 --- GenAI Multi-Modal Issues}
	\begin{block}{Question}
		Multi-modal issues in GenAI validation arise from:
		\begin{enumerate}
			\item Only text data
			\item Processing images, video, and audio in addition to text
			\item Only structured data
			\item Only numerical data
		\end{enumerate}
	\end{block}
	\begin{exampleblock}{Answer: B}\end{exampleblock}
	\begin{block}{Explanation}
		Multi-modal issues: with ability to process images, video, and audio in addition to text, validating GenAI models is more challenging.
	\end{block}
\end{frame}

\begin{frame}{S235 --- GenAI Human Review}
	\begin{block}{Question}
		For GenAI models, human review should be:
		\begin{enumerate}
			\item Avoided
			\item Part of the validation process
			\item Only for deployment
			\item Only for training
		\end{enumerate}
	\end{block}
	\begin{exampleblock}{Answer: B}\end{exampleblock}
	\begin{block}{Explanation}
		Human review must be part of the validation process for GenAI models due to their unique challenges.
	\end{block}
\end{frame}

\begin{frame}{S236 --- GenAI Open-Source Validation}
	\begin{block}{Question}
		Particular attention should be paid to validating:
		\begin{enumerate}
			\item Only in-house models
			\item Open-source and third-party AI models offered as models as a service
			\item Only vendor models
			\item Only traditional models
		\end{enumerate}
	\end{block}
	\begin{exampleblock}{Answer: B}\end{exampleblock}
	\begin{block}{Explanation}
		Attention to open-source and third-party AI models offered as models as a service, which are increasingly common.
	\end{block}
\end{frame}

\begin{frame}{S237 --- GenAI Computational Resources}
	\begin{block}{Question}
		GenAI models require significant computational resources for:
		\begin{enumerate}
			\item Only training
			\item Training, validation, and recalibration
			\item Only inference
			\item Only storage
		\end{enumerate}
	\end{block}
	\begin{exampleblock}{Answer: B}\end{exampleblock}
	\begin{block}{Explanation}
		GenAI models require significant computer resources to train, validate, and recalibrate.
	\end{block}
\end{frame}

\begin{frame}{S238 --- GenAI Interpretability Challenge}
	\begin{block}{Question}
		GenAI model results are difficult to:
		\begin{enumerate}
			\item Generate
			\item Explain or interpret
			\item Store
			\item Delete
		\end{enumerate}
	\end{block}
	\begin{exampleblock}{Answer: B}\end{exampleblock}
	\begin{block}{Explanation}
		GenAI model results are difficult to explain or interpret, making validation challenging.
	\end{block}
\end{frame}

\begin{frame}{S239 --- GenAI Rapid Development}
	\begin{block}{Question}
		The rapidly changing technological development in GenAI requires:
		\begin{enumerate}
			\item Less frequent review
			\item More frequent review and recalibration of models
			\item No review
			\item Only initial validation
		\end{enumerate}
	\end{block}
	\begin{exampleblock}{Answer: B}\end{exampleblock}
	\begin{block}{Explanation}
		GenAI is a rapidly changing area requiring more frequent review and recalibration of models.
	\end{block}
\end{frame}

\begin{frame}{S240 --- GenAI Performance Metrics}
	\begin{block}{Question}
		GenAI performance metrics are:
		\begin{enumerate}
			\item Simpler than traditional ML
			\item More complex and varied than traditional ML
			\item The same as traditional ML
			\item Nonexistent
		\end{enumerate}
	\end{block}
	\begin{exampleblock}{Answer: B}\end{exampleblock}
	\begin{block}{Explanation}
		GenAI performance metrics are more complex and varied than traditional ML due to open-ended, unstructured output.
	\end{block}
\end{frame}

\begin{frame}{S241 --- GenAI Temperature Effect}
	\begin{block}{Question}
		Higher temperature in GenAI models leads to:
		\begin{enumerate}
			\item More predictable output
			\item More creative/random output
			\item No change
			\item Faster output
		\end{enumerate}
	\end{block}
	\begin{exampleblock}{Answer: B}\end{exampleblock}
	\begin{block}{Explanation}
		Higher temperature widens probability distribution, leading to more creative/random output. Lower temperature = more predictable.
	\end{block}
\end{frame}

\begin{frame}{S242 --- GenAI Top-K Effect}
	\begin{block}{Question}
		Lower Top-K in GenAI models leads to:
		\begin{enumerate}
			\item More creativity
			\item Less creativity (more predictable)
			\item No change
			\item More tokens
		\end{enumerate}
	\end{block}
	\begin{exampleblock}{Answer: B}\end{exampleblock}
	\begin{block}{Explanation}
		Lower Top-K restricts to fewer most probable tokens, reducing creativity. Higher Top-K allows more variety.
	\end{block}
\end{frame}

\begin{frame}{S243 --- GenAI Top-P Effect}
	\begin{block}{Question}
		Lower Top-P in GenAI models leads to:
		\begin{enumerate}
			\item More creativity
			\item Less creativity (more predictable)
			\item No change
			\item More tokens
		\end{enumerate}
	\end{block}
	\begin{exampleblock}{Answer: B}\end{exampleblock}
	\begin{block}{Explanation}
		Lower Top-P restricts to smaller set of tokens, reducing creativity. Higher Top-P allows more variety.
	\end{block}
\end{frame}

\begin{frame}{S244 --- GenAI Temperature Range}
	\begin{block}{Question}
		In OpenAI models, temperature ranges from:
		\begin{enumerate}
			\item 0 to 1
			\item 0 to 2
			\item -1 to 1
			\item 1 to 10
		\end{enumerate}
	\end{block}
	\begin{exampleblock}{Answer: B}\end{exampleblock}
	\begin{block}{Explanation}
		OpenAI temperature ranges from 0 to 2. Below 1 = more predictable; above 1 = more random.
	\end{block}
\end{frame}

\begin{frame}{S245 --- GenAI Temperature=1}
	\begin{block}{Question}
		When temperature = 1 in an LLM:
		\begin{enumerate}
			\item Distribution is narrowed
			\item Distribution is widened
			\item Probability distribution is not altered
			\item Model stops working
		\end{enumerate}
	\end{block}
	\begin{exampleblock}{Answer: C}\end{exampleblock}
	\begin{block}{Explanation}
		Temperature = 1 means probability distribution is not altered. Below 1 = narrowed; above 1 = widened/flattened.
	\end{block}
\end{frame}

\begin{frame}{S246 --- GenAI Tokenization}
	\begin{block}{Question}
		Tokenization in GenAI refers to:
		\begin{enumerate}
			\item Removing stop words
			\item Separating text into smaller units (tokens)
			\item Converting words to vectors
			\item Replacing words with synonyms
		\end{enumerate}
	\end{block}
	\begin{exampleblock}{Answer: B}\end{exampleblock}
	\begin{block}{Explanation}
		Tokenization: separating text into smaller units (tokens) for processing by LLMs.
	\end{block}
\end{frame}

\begin{frame}{S247 --- GenAI Embedding}
	\begin{block}{Question}
		Word embeddings in GenAI are:
		\begin{enumerate}
			\item Sparse vectors
			\item Dense vector representations of words
			\item One-hot vectors
			\item Binary vectors
		\end{enumerate}
	\end{block}
	\begin{exampleblock}{Answer: B}\end{exampleblock}
	\begin{block}{Explanation}
		Word embeddings are dense vector representations capturing semantic meaning. Word2Vec and GloVe are examples.
	\end{block}
\end{frame}

\begin{frame}{S248 --- GenAI Attention Mechanism}
	\begin{block}{Question}
		The attention mechanism in transformers:
		\begin{enumerate}
			\item Ignores word relationships
			\item Computes similarity between query and key vectors to weight values
			\item Only processes one word at a time
			\item Uses no weights
		\end{enumerate}
	\end{block}
	\begin{exampleblock}{Answer: B}\end{exampleblock}
	\begin{block}{Explanation}
		Attention: computes similarity between query and key vectors, uses softmax weights to combine value vectors.
	\end{block}
\end{frame}

\begin{frame}{S249 --- GenAI Multi-Head Attention}
	\begin{block}{Question}
		Multi-head attention in transformers:
		\begin{enumerate}
			\item Uses one attention head
			\item Uses multiple attention heads to capture different relationships
			\item Ignores relationships
			\item Uses no weights
		\end{enumerate}
	\end{block}
	\begin{exampleblock}{Answer: B}\end{exampleblock}
	\begin{block}{Explanation}
		Multi-head attention uses several attention heads, each focusing on different aspects of the sequence.
	\end{block}
\end{frame}

\begin{frame}{S250 --- GenAI Positional Embeddings}
	\begin{block}{Question}
		Positional embeddings in transformers:
		\begin{enumerate}
			\item Are not used
			\item Provide information about the position of words in a sequence
			\item Only encode word meaning
			\item Only encode syntax
		\end{enumerate}
	\end{block}
	\begin{exampleblock}{Answer: B}\end{exampleblock}
	\begin{block}{Explanation}
		Positional embeddings provide information about word positions, using sine/cosine functions or learned parameters.
	\end{block}
\end{frame}

\end{document}